\input{preamble}

% D'accord, commençons ici.
%
\begin{document}

\title{Platitude sur les espaces algébriques}

\maketitle

\phantomsection
\label{section-phantom}

\tableofcontents



\section{Introduction}
\label{section-introduction}

\noindent
Dans ce chapitre, nous étudions quelques résultats avancés sur les modules plats et
les morphismes plats dans le cadre des espaces algébriques. Nous recommandons vivement
au lecteur de consulter d'abord le chapitre correspondant
dans le cadre des schémas, voir
Compléments sur la platitude, Section \ref{flat-section-introduction}.
L'article \cite{GruRay} de Raynaud et Gruson constitue une référence.







\section{Impuretés}
\label{section-impure}

\noindent
Cette section est l'analogue de
Compléments sur la platitude, Section \ref{flat-section-impure}.

\begin{situation}
\label{situation-pre-pure}
Soit $S$ un schéma. Soit $f : X \to Y$ un morphisme
de type fini et décent\footnote{Les morphismes quasi-séparés sont décents, voir
Espaces décents, Lemme
\ref{decent-spaces-lemma-properties-trivial-implications}.
Pour tout morphisme $\Spec(k) \to Y$, où $k$ est un corps,
l'espace algébrique $X_k$ est de présentation finie sur $k$
car il est de type fini sur $k$ et quasi-séparé d'après
Espaces décents, Lemme
\ref{decent-spaces-lemma-locally-Noetherian-decent-quasi-separated}.}
entre espaces algébriques sur $S$. De plus, $\mathcal{F}$ est un $\mathcal{O}_X$-module
quasi-cohérent de type fini. Enfin, $y \in |Y|$ est un point de $Y$.
\end{situation}

\noindent
Dans cette situation, considérons un schéma $T$, un morphisme $g : T \to Y$,
un point $t \in T$ tel que $g(t) = y$, une spécialisation $t' \leadsto t$ dans
$T$, et un point $\xi \in |X_T|$ au-dessus de $t'$. Ici $X_T = T \times_Y X$.
On a le diagramme
\begin{equation}
\label{equation-impurity}
\vcenter{
\xymatrix{
\xi \ar@{|->}[d] & \\
t' \ar@{~>}[r] & t \ar@{|->}[r] \ar[r] & y
}
}
\quad\quad
\vcenter{
\xymatrix{
X_T \ar[d]_{f_T} \ar[r] & X \ar[d]^f \\
T \ar[r]^g & Y
}
}
\end{equation}
De plus, notons $\mathcal{F}_T$ l'image réciproque de $\mathcal{F}$ sur $X_T$.

\begin{definition}
\label{definition-impurity}
Dans la
Situation \ref{situation-pre-pure},
nous disons qu'un diagramme (\ref{equation-impurity}) définit une
{\it impureté de $\mathcal{F}$ au-dessus de $y$}
si $\xi \in \text{Ass}_{X_T/T}(\mathcal{F}_T)$ et
$t \not \in f_T(\overline{\{\xi\}})$. Nous l'indiquerons
en disant : « soit $(g : T \to Y, t' \leadsto t, \xi)$
une impureté de $\mathcal{F}$ au-dessus de $y$ ».
\end{definition}

\noindent
Autrement dit, $(g : T \to Y, t' \leadsto t, \xi)$ est
une impureté de $\mathcal{F}$ au-dessus de $y$ s'il n'existe aucune spécialisation
$\xi \leadsto \theta$ dans l'espace topologique $|X_T|$ telle que
$f_T(\theta) = t$. Les spécialisations se comportent mal dans les espaces algébriques
non décents. Si le morphisme $f$ est décent, alors $X_T$
est un espace algébrique décent pour tout morphisme $g : T \to Y$ comme ci-dessus, voir
Espaces décents, Définition \ref{decent-spaces-definition-relative-conditions}.

\begin{lemma}
\label{lemma-impure-limit}
Dans la Situation \ref{situation-pre-pure}.
Soit $(g : T \to S, t' \leadsto t, \xi)$ une impureté de
$\mathcal{F}$ au-dessus de $y$. Supposons que $T = \lim_{i \in I} T_i$ soit une limite dirigée
de schémas affines sur $Y$. Alors, pour un certain $i$, le triplet
$(T_i \to Y, t'_i \leadsto t_i, \xi_i)$ est une impureté de
$\mathcal{F}$ au-dessus de $y$.
\end{lemma}

\begin{proof}
Les notations de l'énoncé signifient ceci : soient $p_i : T \to T_i$
les morphismes de projection, et posons $t_i = p_i(t)$ et $t'_i = p_i(t')$.
Enfin, $\xi_i \in |X_{T_i}|$ est l'image de $\xi$. D'après
Diviseurs sur les espaces, Lemme
\ref{spaces-divisors-lemma-base-change-relative-assassin}
nous avons $\xi_i \in \text{Ass}_{X_{T_i}/T_i}(\mathcal{F}_{T_i})$.
Il reste donc seulement à montrer que
$t_i \not \in f_{T_i}(\overline{\{\xi_i\}})$ pour un certain $i$.

\medskip\noindent
Soit $Z_i \subset X_{T_i}$ la structure de schéma réduite induite
sur $\overline{\{\xi_i\}} \subset |X_{T_i}|$,
et soit $Z \subset X_T$ la structure de schéma réduite induite sur
$\overline{\{\xi\}} \subset |X_T|$.
Alors $Z = \lim Z_i$ d'après
Limites d'espaces, Lemme \ref{spaces-limits-lemma-inverse-limit-irreducibles}
(le lemme s'applique car chaque $X_{T_i}$ est décent).
Choisissons un corps $k$ et un morphisme $\Spec(k) \to T$ dont l'image est $t$.
Alors
$$
\emptyset =
Z \times_T \Spec(k) = (\lim Z_i) \times_{(\lim T_i)} \Spec(k)
= \lim Z_i \times_{T_i} \Spec(k)
$$
car les limites commutent aux produits fibrés (les limites commutent aux limites).
Chaque $Z_i \times_{T_i} \Spec(k)$ est quasi-compact, car $X_{T_i} \to T_i$
est de type fini, et donc $Z_i \to T_i$ est de type fini.
Ainsi $Z_i \times_{T_i} \Spec(k)$ est vide pour un certain $i$ d'après
Limites d'espaces, Lemme \ref{spaces-limits-lemma-limit-nonempty}.
Puisque l'image du composé $\Spec(k) \to T \to T_i$ est $t_i$,
nous obtenons le résultat voulu.
\end{proof}

\noindent
Les impuretés remontent par changement de base plat.

\begin{lemma}
\label{lemma-flat-ascent-impurity}
Dans la Situation \ref{situation-pre-pure}.
Soit $(Y_1, y_1) \to (Y, y)$ un morphisme d'espaces algébriques
pointés sur $S$. Supposons que $Y_1 \to Y$ soit plat en $y_1$.
Si $(T \to Y, t' \leadsto t, \xi)$ est une impureté de
$\mathcal{F}$ au-dessus de $y$, alors il existe une impureté
$(T_1 \to Y_1, t_1' \leadsto t_1, \xi_1)$ de l'image réciproque
$\mathcal{F}_1$ de $\mathcal{F}$ sur $X_1 = Y_1 \times_Y X$
au-dessus de $y_1$, telle que $T_1$ soit \'etale sur $Y_1 \times_Y T$.
\end{lemma}

\begin{proof}
Choisissons un morphisme \'etale $T_1 \to Y_1 \times_Y T$, où $T_1$
est un schéma, et un point $t_1 \in T_1$ s'envoyant sur $y_1$ et $t$.
L'existence d'un tel couple $(T_1, t_1)$ résulte de
Propriétés des espaces, Lemme \ref{spaces-properties-lemma-points-cartesian}.
Le morphisme de schémas $T_1 \to T$ est plat en $t_1$
(utiliser Morphismes d'espaces, Lemme \ref{spaces-morphisms-lemma-base-change-flat}
et la définition des morphismes plats d'espaces algébriques) ;
il existe donc une spécialisation $t'_1 \leadsto t_1$ au-dessus de
$t' \leadsto t$, voir
Morphismes, Lemme \ref{morphisms-lemma-generalizations-lift-flat}.
Choisissons un point $\xi_1 \in |X_{T_1}|$ s'envoyant sur $t'_1$
et sur $\xi$, avec $\xi_1 \in \text{Ass}_{X_{T_1}/T_1}(\mathcal{F}_{T_1})$.
Autrement dit, un point de $\Spec(\kappa(t'_1) \otimes_{\kappa(t')} \kappa(\xi))$.
C'est possible d'après
Diviseurs sur les espaces, Lemme
\ref{spaces-divisors-lemma-base-change-relative-assassin}.
Comme l'adhérence $Z_1$ de $\{\xi_1\}$ dans $|X_{T_1}|$ s'envoie dans
l'adhérence de $\{\xi\}$ dans $|X_T|$, nous en concluons que
l'image de $Z_1$ dans $|T_1|$ ne peut contenir $t_1$.
Ainsi $(T_1 \to Y_1, t'_1 \leadsto t_1, \xi_1)$
est une impureté de $\mathcal{F}_1$ au-dessus de $Y_1$.
\end{proof}

\begin{lemma}
\label{lemma-pure-along-X-y}
Dans la Situation \ref{situation-pre-pure}. Soit $\overline{y}$ un point
géométrique au-dessus de $y$. Soit $\mathcal{O} = \mathcal{O}_{Y, \overline{y}}$
l'anneau local étale de $Y$ en $\overline{y}$. Notons
$Y^{sh} = \Spec(\mathcal{O})$, $X^{sh} = X \times_Y Y^{sh}$, et
$\mathcal{F}^{sh}$ l'image réciproque de $\mathcal{F}$ sur $X^{sh}$.
Les conditions suivantes sont équivalentes :
\begin{enumerate}
\item il existe une impureté
$(Y^{sh} \to Y, y' \leadsto \overline{y}, \xi)$
de $\mathcal{F}$ au-dessus de $y$,
\item tout point de $\text{Ass}_{X^{sh}/Y^{sh}}(\mathcal{F}^{sh})$
se spécialise en un point de la fibre fermée $X_{\overline{y}}$,
\item il existe une impureté $(T \to Y, t' \leadsto t, \xi)$
de $\mathcal{F}$ au-dessus de $y$ telle que $(T, t) \to (Y, y)$ soit un
voisinage \'etale, et
\item il existe une impureté $(T \to Y, t' \leadsto t, \xi)$
de $\mathcal{F}$ au-dessus de $y$ telle que $T \to Y$ soit quasi-fini en $t$.
\end{enumerate}
\end{lemma}

\begin{proof}
L'équivalence de (1) et (2) résulte immédiatement de la définition.

\medskip\noindent
Rappelons que $\mathcal{O} = \mathcal{O}_{Y, \overline{y}}$
est la colimite filtrante des $\mathcal{O}(V)$ sur la catégorie
des voisinages \'etales $(V, \overline{v}) \to (Y, \overline{y})$
(Propriétés des espaces, Lemme \ref{spaces-properties-lemma-cofinal-etale}).
De plus, il suffit de considérer les voisinages \'etales affines $V$.
Par conséquent, $Y^{sh} = \Spec(\mathcal{O}) = \lim \Spec(\mathcal{O}(V)) = \lim V$.
Le Lemme \ref{lemma-impure-limit} montre donc que (1) implique (3).

\medskip\noindent
Puisqu'un morphisme \'etale est localement quasi-fini
(Morphismes d'espaces, Lemme
\ref{spaces-morphisms-lemma-etale-locally-quasi-finite})
nous voyons que (3) implique (4).

\medskip\noindent
Enfin, supposons (4). Après avoir remplacé $T$ par un voisinage ouvert de $t$,
nous pouvons supposer que $T \to Y$ est localement quasi-fini.
Le Lemme \ref{lemma-flat-ascent-impurity}
fournit une impureté
$(T_1 \to Y^{sh}, t_1' \leadsto t_1, \xi_1)$
pour laquelle $T_1 \to T \times_Y Y^{sh}$ est
\'etale. Puisqu'un morphisme \'etale est localement quasi-fini,
Morphismes d'espaces, Lemme
\ref{spaces-morphisms-lemma-base-change-quasi-finite}, et
Morphismes, Lemme \ref{morphisms-lemma-composition-quasi-finite},
montrent que $T_1 \to Y^{sh}$ est localement quasi-fini.
Comme $\mathcal{O}$ est strictement hensélien, nous pouvons appliquer Compléments sur les morphismes, Lemme
\ref{more-morphisms-lemma-etale-makes-quasi-finite-finite-at-point}
pour voir qu'après avoir remplacé $T_1$ par un voisinage ouvert et fermé
de $t_1$, nous pouvons supposer que $T_1 \to Y^{sh} = \Spec(\mathcal{O})$
est fini. Soit $\theta \in |X^{sh}|$ l'image de
$\xi_1$, et soit $y' \in \Spec(\mathcal{O})$ l'image
de $t_1'$. D'après Diviseurs sur les espaces, Lemme
\ref{spaces-divisors-lemma-base-change-relative-assassin}
nous voyons que $\theta \in \text{Ass}_{X^{sh}/Y^{sh}}(\mathcal{F}^{sh})$.
Comme $\pi : X_{T_1} \to X^{sh}$
est fini, il induit une application fermée $|X_{T_1}| \to |X^{sh}|$.
L'image de $\overline{\{\xi_1\}}$ est donc $\overline{\{\theta\}}$.
Il s'ensuit que $(Y^{sh} \to Y, y' \leadsto \overline{y}, \theta)$
est une impureté de $\mathcal{F}$ au-dessus de $y$, ce qui achève la démonstration.
\end{proof}








\section{Modules relativement purs}
\label{section-pure}

\noindent
Cette section est l'analogue de
Compléments sur la platitude, Section \ref{flat-section-pure}.

\begin{definition}
\label{definition-pure}
Dans la Situation \ref{situation-pre-pure}.
\begin{enumerate}
\item Nous disons que $\mathcal{F}$ est {\it pur au-dessus de $y$} si {\bf aucune} des
conditions équivalentes du Lemme \ref{lemma-pure-along-X-y} n'est satisfaite.
\item Nous disons que $\mathcal{F}$ est {\it universellement pur au-dessus de $y$}
s'il n'existe aucune impureté de $\mathcal{F}$ au-dessus de $y$.
\item Nous disons que $X$ est {\it pur au-dessus de $y$} si $\mathcal{O}_X$
est pur au-dessus de $y$.
\item Nous disons que $\mathcal{F}$ est {\it universellement $Y$-pur}, ou
{\it universellement pur relativement à $Y$}, si $\mathcal{F}$ est universellement
pur au-dessus de $y$ pour tout $y \in |Y|$.
\item Nous disons que $\mathcal{F}$ est {\it $Y$-pur}, ou
{\it pur relativement à $Y$}, si $\mathcal{F}$ est pur au-dessus de $y$
pour tout $y \in |Y|$.
\item Nous disons que $X$ est {\it $Y$-pur} ou {\it pur relativement à $Y$}
si $\mathcal{O}_X$ est pur relativement à $Y$.
\end{enumerate}
\end{definition}

\noindent
Voici les lemmes obligés.

\begin{lemma}
\label{lemma-base-change-universally}
Dans la Situation \ref{situation-pre-pure}.
\begin{enumerate}
\item $\mathcal{F}$ est universellement pur au-dessus de $y$, et
\item pour tout morphisme $(Y', y') \to (Y, y)$ d'espaces algébriques pointés,
l'image réciproque $\mathcal{F}_{Y'}$ est pure au-dessus de $y'$.
\end{enumerate}
En particulier, $\mathcal{F}$ est universellement pur relativement à $Y$ si et
seulement si tout changement de base $\mathcal{F}_{Y'}$ de $\mathcal{F}$ est
pur relativement à $Y'$.
\end{lemma}

\begin{proof}
C'est formel.
\end{proof}

\begin{lemma}
\label{lemma-quasi-finite-base-change}
Dans la Situation \ref{situation-pre-pure}.
Soit $(Y', y') \to (Y, y)$ un morphisme d'espaces algébriques pointés.
Si $Y' \to Y$ est quasi-fini en $y'$ et si $\mathcal{F}$ est pur au-dessus de $y$,
alors $\mathcal{F}_{Y'}$ est pur au-dessus de $y'$.
\end{lemma}

\begin{proof}
Si $(T \to Y', t' \leadsto t, \xi)$ est une impureté de
$\mathcal{F}_{Y'}$ au-dessus de $y'$, avec $T \to Y'$ quasi-fini en $t$,
alors $(T \to Y, t' \to t, \xi)$ est une impureté de $\mathcal{F}$
au-dessus de $y$, avec $T \to Y$ quasi-fini en $t$ ; voir
Morphismes d'espaces, Lemme
\ref{spaces-morphisms-lemma-composition-quasi-finite}.
Le lemme résulte donc immédiatement de la définition de la pureté.
\end{proof}

\noindent
La pureté satisfait la descente plate.

\begin{lemma}
\label{lemma-flat-descend-pure}
Dans la Situation \ref{situation-pre-pure}.
Soit $(Y_1, y_1) \to (Y, y)$ un morphisme d'espaces algébriques pointés.
Supposons que $Y_1 \to Y$ soit plat en $y_1$.
\begin{enumerate}
\item Si $\mathcal{F}_{Y_1}$ est pur au-dessus de $y_1$,
alors $\mathcal{F}$ est pur au-dessus de $y$.
\item Si $\mathcal{F}_{Y_1}$ est universellement pur au-dessus de $y_1$,
alors $\mathcal{F}$ est universellement pur au-dessus de $y$.
\end{enumerate}
\end{lemma}

\begin{proof}
Cela vient de ce que les impuretés remontent par changement de base plat, voir
Lemme \ref{lemma-flat-ascent-impurity}. Par exemple,
pour (1), toute impureté $(T \to Y, t' \leadsto t, \xi)$
de $\mathcal{F}$ au-dessus de $y$, avec $T \to Y$ quasi-fini en $t$,
conduit par ce lemme à une impureté
$(T_1 \to Y_1, t_1' \leadsto t_1, \xi_1)$ de l'image réciproque
$\mathcal{F}_1$ de $\mathcal{F}$ sur $X_1 = Y_1 \times_Y X$
au-dessus de $y_1$, telle que $T_1$ soit \'etale sur $Y_1 \times_Y T$.
Ainsi $T_1 \to Y_1$ est quasi-fini en $t_1$, car
les morphismes \'etales sont localement quasi-finis et les composés
de morphismes localement quasi-finis sont localement quasi-finis
(Morphismes d'espaces, Lemmes
\ref{spaces-morphisms-lemma-etale-locally-quasi-finite} et
\ref{spaces-morphisms-lemma-composition-quasi-finite}).
Le raisonnement est analogue pour (2).
\end{proof}

\begin{lemma}
\label{lemma-supported-on-closed}
Dans la Situation \ref{situation-pre-pure}. Soit $i : Z \to X$ une immersion fermée,
et supposons que $\mathcal{F} = i_*\mathcal{G}$ pour un faisceau
quasi-cohérent de type fini $\mathcal{G}$ sur $Z$.
Alors $\mathcal{G}$ est (universellement) pur au-dessus de $y$
si et seulement si $\mathcal{F}$ est (universellement) pur au-dessus de $y$.
\end{lemma}

\begin{proof}
Cela résulte de Diviseurs sur les espaces, Lemme
\ref{spaces-divisors-lemma-relative-weak-assassin-finite}.
\end{proof}

\begin{lemma}
\label{lemma-proper-pure}
Dans la Situation \ref{situation-pre-pure}.
\begin{enumerate}
\item Si le support de $\mathcal{F}$ est propre sur $Y$, alors
$\mathcal{F}$ est universellement pur relativement à $Y$.
\item Si $f$ est propre, alors
$\mathcal{F}$ est universellement pur relativement à $Y$.
\item Si $f$ est propre, alors $X$ est universellement pur relativement à $Y$.
\end{enumerate}
\end{lemma}

\begin{proof}
Réduisons d'abord (1) à (2). Soit en effet $Z \subset X$ le
support schématique de $\mathcal{F}$
(Morphismes d'espaces, Définition
\ref{spaces-morphisms-definition-scheme-theoretic-support}). Soit $i : Z \to X$
l'immersion fermée correspondante, et écrivons
$\mathcal{F} = i_*\mathcal{G}$ pour un $\mathcal{O}_Z$-module $\mathcal{G}$
quasi-cohérent de type fini.
Dans le cas (1), $Z \to Y$ est propre par hypothèse.
Le Lemme \ref{lemma-supported-on-closed} réduit donc le cas (1) au cas (2).

\medskip\noindent
Supposons $f$ propre.
Soit $(g : T \to Y, t' \leadsto t, \xi)$ une impureté de $\mathcal{F}$
au-dessus de $y$. Puisque $f$ est propre, il est universellement fermé. Par conséquent,
$f_T : X_T \to T$ est fermé. Comme $f_T(\xi) = t'$, cela implique que
$t \in f(\overline{\{\xi\}})$, contradiction.
\end{proof}







\section{Modules plats de type fini}
\label{section-finite-type-flat}

\noindent
Comparer avec
Compléments sur la platitude, Sections
\ref{flat-section-finite-type-flat-I},
\ref{flat-section-finite-type-flat-II}, et
\ref{flat-section-finite-type-flat-III}.
La plupart de ces résultats ont des conséquences immédiates
pour les espaces algébriques par localisation \'etale.

\begin{lemma}
\label{lemma-existence-complete}
Soit $S$ un schéma.
Soit $X \to Y$ un morphisme de type fini d'espaces algébriques sur $S$.
Soit $\mathcal{F}$ un $\mathcal{O}_X$-module quasi-cohérent de type fini.
Soit $y \in |Y|$ un point. Il existe un morphisme \'etale
$(Y', y') \to (Y, y)$, avec $Y'$ schéma affine, et des morphismes \'etales
$h_i : W_i \to X_{Y'}$, $i = 1, \ldots, n$, tels que, pour tout
$i$, il existe un dévissage complet de $\mathcal{F}_i/W_i/Y'$ au-dessus de $y'$,
où $\mathcal{F}_i$ est l'image réciproque de $\mathcal{F}$ sur $W_i$,
et tels que $|(X_{Y'})_{y'}| \subset \bigcup h_i(W_i)$.
\end{lemma}

\begin{proof}
La question est locale pour la topologie \'etale sur $Y$ ; nous pouvons donc supposer que $Y$
est un schéma affine. Alors $X$ est quasi-compact, de sorte que nous pouvons
choisir un schéma affine $X'$ et un morphisme \'etale surjectif
$X' \to X$. Nous pouvons alors appliquer
Compléments sur la platitude, Lemme \ref{flat-lemma-existence-complete}
à $X' \to Y$, $(X' \to Y)^*\mathcal{F}$ et $y$ pour
obtenir le résultat voulu.
\end{proof}

\begin{lemma}
\label{lemma-open-in-fibre-where-flat}
Soit $S$ un schéma.
Soit $f : X \to Y$ un morphisme d'espaces algébriques sur $S$
qui est localement de type fini.
Soit $\mathcal{F}$ un $\mathcal{O}_X$-module quasi-cohérent de type fini.
Soit $y \in |Y|$ et $F = f^{-1}(\{y\}) \subset |X|$. Alors l'ensemble
$$
\{x \in F \mid \mathcal{F} \text{ est plat sur }Y\text{ en }x\}
$$
est ouvert dans $F$.
\end{lemma}

\begin{proof}
Choisissons un schéma $V$, un point $v \in V$ et un morphisme \'etale $V \to Y$
envoyant $v$ sur $y$. Choisissons un schéma $U$ et un morphisme \'etale surjectif
$U \to V \times_Y X$. Alors $|U_v| \to F$ est une application continue ouverte
d'espaces topologiques, puisque $|U| \to |X|$ est continue et ouverte.
Le résultat découle donc du cas des schémas, à savoir
Compléments sur la platitude, Lemme \ref{flat-lemma-open-in-fibre-where-flat}.
\end{proof}

\begin{lemma}
\label{lemma-bourbaki-finite-type-general-base-at-point}
Soit $S$ un schéma.
Soit $f : X \to Y$ un morphisme d'espaces algébriques sur $S$ qui est
localement de type fini. Soit $x \in |X|$, d'image $y \in |Y|$.
Soit $\mathcal{F}$ un faisceau quasi-cohérent de type fini sur $X$.
Soit $\mathcal{G}$ un faisceau quasi-cohérent sur $Y$.
Si $\mathcal{F}$ est plat en $x$ sur $Y$, alors
$$
x \in \text{WeakAss}_X(\mathcal{F} \otimes_{\mathcal{O}_X} f^*\mathcal{G})
\Leftrightarrow
y \in \text{WeakAss}_Y(\mathcal{G})
\text{ et }
x \in \text{Ass}_{X/Y}(\mathcal{F}).
$$
\end{lemma}

\begin{proof}
Choisissons un diagramme commutatif
$$
\xymatrix{
U \ar[d] \ar[r]_g & V \ar[d] \\
X \ar[r]^f & Y
}
$$
où $U$ et $V$ sont des schémas et les flèches verticales sont \'etales surjectives.
Choisissons $u \in U$ s'envoyant sur $x$. Posons $\mathcal{E} = \mathcal{F}|_U$
et $\mathcal{H} = \mathcal{G}|_V$.
Soit $v \in V$ l'image de $u$. Alors
$x \in \text{WeakAss}_X(\mathcal{F} \otimes_{\mathcal{O}_X} f^*\mathcal{G})$
si et seulement si
$u \in \text{WeakAss}_X(\mathcal{E} \otimes_{\mathcal{O}_X} g^*\mathcal{H})$
d'après Diviseurs sur les espaces, Définition
\ref{spaces-divisors-definition-weakly-associated}.
De même, $y \in \text{WeakAss}_Y(\mathcal{G})$ si et seulement si
$v \in \text{WeakAss}_V(\mathcal{H})$.
Enfin, $x \in \text{Ass}_{X/Y}(\mathcal{F})$ si et seulement si
$u \in \text{Ass}_{U_v}(\mathcal{E}|_{U_v})$ d'après
Diviseurs sur les espaces, Définition
\ref{spaces-divisors-definition-relative-weak-assassin}.
Observons que la platitude de $\mathcal{F}$ en $x$ est
équivalente à celle de $\mathcal{E}$ en $u$, voir
Morphismes d'espaces, Définition \ref{spaces-morphisms-definition-flat-module}.
L'équivalence pour $g : U \to V$, $\mathcal{E}$, $\mathcal{H}$, $u$ et $v$
est Compléments sur la platitude, Lemme
\ref{flat-lemma-bourbaki-finite-type-general-base-at-point}.
\end{proof}

\begin{lemma}
\label{lemma-bourbaki-finite-type-general-base}
Soit $S$ un schéma. Soit $f : X \to Y$ un morphisme d'espaces algébriques
sur $S$ qui est localement de type fini.
Soit $\mathcal{F}$ un faisceau quasi-cohérent de type fini sur $X$
qui est plat sur $Y$. Soit $\mathcal{G}$ un faisceau quasi-cohérent sur $Y$.
Alors
$$
\text{WeakAss}_X(\mathcal{F} \otimes_{\mathcal{O}_X} f^*\mathcal{G}) =
\text{Ass}_{X/Y}(\mathcal{F}) \cap
|f|^{-1}(\text{WeakAss}_Y(\mathcal{G}))
$$
\end{lemma}

\begin{proof}
C'est une conséquence immédiate du
Lemme \ref{lemma-bourbaki-finite-type-general-base-at-point}.
\end{proof}

\begin{theorem}
\label{theorem-finite-type-flat}
Soit $S$ un schéma.
Soit $f : X \to Y$ un morphisme d'espaces algébriques sur $S$.
Soit $\mathcal{F}$ un $\mathcal{O}_X$-module quasi-cohérent.
Supposons que
\begin{enumerate}
\item $X \to Y$ soit localement de présentation finie,
\item $\mathcal{F}$ soit un $\mathcal{O}_X$-module de type fini, et
\item l'ensemble des points faiblement associés de $Y$ soit localement fini dans $Y$.
\end{enumerate}
Alors $U = \{x \in |X| : \mathcal{F}\text{ est plat en }x\text{ sur }Y\}$
est ouvert dans $X$ et $\mathcal{F}|_U$ est un $\mathcal{O}_U$-module
de présentation finie et plat sur $Y$.
\end{theorem}

\begin{proof}
La condition (3) signifie que, si $V \to Y$ est un morphisme \'etale surjectif,
où $V$ est un schéma, alors les points faiblement associés de $V$
sont localement en nombre fini sur le schéma $V$. (Rappelons que les points faiblement
associés de $V$ sont exactement l'image réciproque des points faiblement associés
de $Y$ d'après Diviseurs sur les espaces, Définition
\ref{spaces-divisors-definition-weakly-associated}.)
Cela étant, la question
est locale pour la topologie \'etale sur $X$ et $Y$ ; nous pouvons donc supposer que $X$ et $Y$
sont des schémas. Le résultat découle alors de
Compléments sur la platitude, Théorème \ref{flat-theorem-finite-type-flat}.
\end{proof}

\begin{lemma}
\label{lemma-finite-type-flat-along-fibre-free-variant}
Soit $S$ un schéma. Soit $f : X \to Y$ un morphisme d'espaces algébriques
sur $S$. Soit $\mathcal{F}$ un faisceau quasi-cohérent sur $X$.
Soit $y \in |Y|$. Posons $F = f^{-1}(\{y\}) \subset |X|$. Supposons que
\begin{enumerate}
\item $f$ soit de type fini,
\item $\mathcal{F}$ soit de type fini, et
\item $\mathcal{F}$ soit plat sur $Y$ en tout $x \in F$.
\end{enumerate}
Alors il existe un morphisme \'etale $(Y', y') \to (Y, y)$,
où $Y'$ est un schéma, et un diagramme commutatif d'espaces algébriques
$$
\xymatrix{
X \ar[d] & X' \ar[l]^g \ar[d] \\
Y & \Spec(\mathcal{O}_{Y', y'}) \ar[l]
}
$$
tels que $X' \to X \times_Y \Spec(\mathcal{O}_{Y', y'})$
soit \'etale, que $|X'_{y'}| \to F$ soit surjectif, que $X'$ soit affine,
et que $\Gamma(X', g^*\mathcal{F})$ soit un $\mathcal{O}_{Y', y'}$-module libre.
\end{lemma}

\begin{proof}
Choisissons un morphisme \'etale $(Y', y') \to (Y, y)$, où $Y'$ est un
schéma affine. Alors $X \times_Y Y'$ est quasi-compact.
Choisissons un schéma affine $X'$ et un morphisme \'etale surjectif
$X' \to X \times_Y Y'$. On a le diagramme
$$
\xymatrix{
X \ar[d] & X' \ar[l]^g \ar[d] \\
Y & Y' \ar[l]
}
$$
Alors $\mathcal{F}' = g^*\mathcal{F}$ est plat sur $Y'$ en tout
point de $X'_{y'}$, voir Morphismes d'espaces, Lemme
\ref{spaces-morphisms-lemma-base-change-module-flat}.
Nous pouvons donc appliquer le lemme dans le cas des schémas
(Compléments sur la platitude, Lemme
\ref{flat-lemma-finite-type-flat-along-fibre-free-variant})
au morphisme
$X' \to Y'$, au faisceau quasi-cohérent $g^*\mathcal{F}$ et au point $y'$.
On obtient un morphisme \'etale $(Y'', y'') \to (Y', y')$ et un
diagramme commutatif
$$
\xymatrix{
X \ar[d] & X' \ar[l]^g \ar[d] & X'' \ar[l]^{g'} \ar[d] \\
Y & Y' \ar[l] & \Spec(\mathcal{O}_{Y'', y''}) \ar[l]
}
$$
Pour obtenir le résultat voulu, on prend $(Y'', y'') \to (Y, y)$
et $g \circ g' : X'' \to X$.
\end{proof}

\begin{theorem}
\label{theorem-check-flatness-at-associated-points}
Soit $S$ un schéma.
Soit $f : X \to Y$ un morphisme d'espaces algébriques sur $S$
qui est localement de type fini.
Soit $\mathcal{F}$ un $\mathcal{O}_X$-module quasi-cohérent de type fini.
Soit $x \in |X|$, d'image $y \in |Y|$.
Posons $F = f^{-1}(\{y\}) \subset |X|$.
Considérons les conditions suivantes :
\begin{enumerate}
\item $\mathcal{F}$ est plat en $x$ sur $Y$, et
\item pour tout $x' \in F \cap \text{Ass}_{X/Y}(\mathcal{F})$ qui
se spécialise en $x$, le module $\mathcal{F}$ est plat en $x'$ sur $Y$.
\end{enumerate}
On a toujours (2) $\Rightarrow$ (1). Si $X$ et $Y$ sont
décents, alors (1) $\Rightarrow$ (2).
\end{theorem}

\begin{proof}
Supposons (2).
Choisissons un schéma $V$ et un morphisme \'etale surjectif $V \to Y$.
Choisissons un schéma $U$ et un morphisme \'etale surjectif $U \to V \times_Y X$.
Choisissons un point $u \in U$ s'envoyant sur $x$. Soit $v \in V$ l'image de $u$.
Nous déduirons le résultat du résultat correspondant pour
$\mathcal{F}|_U = (U \to X)^*\mathcal{F}$ et le point $u$.
$U_v$. Cela fonctionne car $\text{Ass}_{U/V}(\mathcal{F}|_U) \cap |U_v|$
est égal à $\text{Ass}_{U_v}(\mathcal{F}|_{U_v})$, ainsi qu'à l'image
réciproque de $F \cap \text{Ass}_{X/Y}(\mathcal{F})$.
Comme l'application $|U_v| \to F$ est continue, nous voyons que
les spécialisations dans $|U_v|$ s'envoient sur des spécialisations dans $F$ ;
la condition (2) est donc héritée par $U \to V$,
$\mathcal{F}|_U$ et le point $u$.
Ainsi Compléments sur la platitude, Théorème
\ref{flat-theorem-check-flatness-at-associated-points} s'applique
et nous concluons que (1) est satisfaite.

\medskip\noindent
Si $Y$ est décent, nous pouvons représenter
$y$ par un monomorphisme quasi-compact $\Spec(k) \to Y$
(par définition des espaces décents, voir
Espaces décents, Définition \ref{decent-spaces-definition-very-reasonable}).
Alors $F = |X_k|$, voir
Espaces décents, Lemme \ref{decent-spaces-lemma-topology-fibre}.
Si, de plus, $X$ est décent (ou, plus généralement, si $f$ est décent, voir
Espaces décents, Définition \ref{decent-spaces-definition-relative-conditions}
et Espaces décents, Lemme
\ref{decent-spaces-lemma-property-for-morphism-out-of-property}),
alors $X_y$ est lui aussi un espace décent. De plus, les spécialisations dans
$F$ se relèvent en spécialisations
dans $U_v \to X_y$, voir
Espaces décents, Lemme \ref{decent-spaces-lemma-decent-specialization}.
Cela étant, il est clair que l'implication réciproque
est vraie, puisqu'elle l'est dans le cas des schémas.
\end{proof}

\begin{lemma}
\label{lemma-check-along-closed-fibre}
Soit $S$ un schéma local de point fermé $s$.
Soit $f : X \to S$ un morphisme d'un espace algébrique $X$ vers $S$
qui est localement de type fini.
Soit $\mathcal{F}$ un $\mathcal{O}_X$-module quasi-cohérent de type fini.
Supposons que
\begin{enumerate}
\item tout point de $\text{Ass}_{X/S}(\mathcal{F})$ se spécialise
en un point de la fibre fermée $X_s$\footnote{C'est par exemple le cas si
$f$ est de type fini et $\mathcal{F}$ est pur le long de $X_s$, ou
si $f$ est propre.},
\item $\mathcal{F}$ est plat sur $S$ en tout point de $X_s$.
\end{enumerate}
Alors $\mathcal{F}$ est plat sur $S$.
\end{lemma}

\begin{proof}
Cela résulte immédiatement du fait qu'il suffit de vérifier la
platitude aux points de l'assassin relatif de $\mathcal{F}$
sur $S$, d'après le
Théorème \ref{theorem-check-flatness-at-associated-points}.
\end{proof}




\section{Modules plats de présentation finie}
\label{section-finitely-presented-flat}

\noindent
Cette section est l'analogue de Compléments sur la platitude, Section
\ref{flat-section-finitely-presented-flat}.

\begin{proposition}
\label{proposition-finite-presentation-flat-at-point}
Soit $S$ un schéma.
Soit $f : X \to Y$ un morphisme d'espaces algébriques sur $S$.
Soit $\mathcal{F}$ un faisceau quasi-cohérent sur $X$.
Soit $x \in |X|$, d'image $y \in |Y|$.
Supposons que
\begin{enumerate}
\item $f$ soit localement de présentation finie,
\item $\mathcal{F}$ soit de présentation finie, et
\item $\mathcal{F}$ soit plat en $x$ sur $Y$.
\end{enumerate}
Alors il existe un diagramme commutatif de schémas pointés
$$
\xymatrix{
(X, x) \ar[d] & (X', x') \ar[l]^g \ar[d] \\
(Y, y) & (Y', y') \ar[l]
}
$$
dont les flèches horizontales sont \'etales, tel que $X'$ et $Y'$
soient affines et que
$\Gamma(X', g^*\mathcal{F})$ soit un
$\Gamma(Y', \mathcal{O}_{Y'})$-module projectif.
\end{proposition}

\begin{proof}
Sous cette forme, la proposition se réduit immédiatement
au cas des schémas, qui est
Compléments sur la platitude, Proposition
\ref{flat-proposition-finite-presentation-flat-at-point}.
\end{proof}

\begin{lemma}
\label{lemma-finite-presentation-flat-along-fibre}
Soit $S$ un schéma.
Soit $f : X \to Y$ un morphisme d'espaces algébriques sur $S$.
Soit $\mathcal{F}$ un faisceau quasi-cohérent sur $X$.
Soit $y \in |Y|$. Posons $F = f^{-1}(\{y\}) \subset |X|$. Supposons que
\begin{enumerate}
\item $f$ soit de présentation finie,
\item $\mathcal{F}$ soit de présentation finie, et
\item $\mathcal{F}$ soit plat sur $Y$ en tout $x \in F$.
\end{enumerate}
Alors il existe un diagramme commutatif d'espaces algébriques
$$
\xymatrix{
X \ar[d] & X' \ar[l]^g \ar[d] \\
Y & Y' \ar[l]_h
}
$$
tel que $h$ et $g$ soient \'etales, qu'il existe un point
$y' \in |Y'|$ s'envoyant sur $y$, que $F \subset g(|X'|)$,
que les espaces algébriques $X'$ et $Y'$ soient affines, et que
$\Gamma(X', g^*\mathcal{F})$ soit un
$\Gamma(Y', \mathcal{O}_{Y'})$-module projectif.
\end{lemma}

\begin{proof}
Sous cette forme, le lemme se réduit immédiatement
au cas des schémas, qui est
Compléments sur la platitude, Lemme
\ref{flat-lemma-finite-presentation-flat-along-fibre}.
\end{proof}




\section{Un critère de pureté}
\label{section-criterion-purity}

\noindent
Cette section est l'analogue de Compléments sur la platitude, Section
\ref{flat-section-criterion-purity}.

\begin{lemma}
\label{lemma-associated-point-specializes}
Soit $S$ un schéma. Soit $X$ un espace algébrique décent
localement de type fini sur $S$.
Soit $\mathcal{F}$ un $\mathcal{O}_X$-module quasi-cohérent de type fini.
Soit $s \in S$ tel que $\mathcal{F}$ soit plat sur $S$ en tout point de $X_s$.
Soit $x' \in \text{Ass}_{X/S}(\mathcal{F})$. Si
l'adhérence de $\{x'\}$ dans $|X|$ rencontre $|X_s|$, alors elle
rencontre $\text{Ass}_{X/S}(\mathcal{F}) \cap |X_s|$.
\end{lemma}

\begin{proof}
Observons que $|X_s| \subset |X|$ est l'ensemble des points de $|X|$
au-dessus de $s \in S$, voir
Espaces décents, Lemme \ref{decent-spaces-lemma-topology-fibre}.
Soit $t \in |X_s|$ une spécialisation de $x'$ dans $|X|$.
Choisissons un schéma affine $U$, un point $u \in U$ et
un morphisme \'etale $\varphi : U \to X$ envoyant $u$ sur $t$.
D'après Espaces décents, Lemme \ref{decent-spaces-lemma-decent-specialization},
nous pouvons choisir une spécialisation $u' \leadsto u$
telle que $u'$ s'envoie sur $x'$. Posons $g = f \circ \varphi$.
Observons que $s' = g(u') = f(x')$ se spécialise en $s$.
Par notre définition de $\text{Ass}_{X/S}(\mathcal{F})$,
nous avons $u' \in \text{Ass}_{U/S}(\varphi^*\mathcal{F})$.
La version pour les schémas de ce lemme
(Compléments sur la platitude, Lemme \ref{flat-lemma-associated-point-specializes})
montre qu'il existe une spécialisation $u' \leadsto u$ avec
$u \in \text{Ass}_{U_s}(\varphi^*\mathcal{F}_s) =
\text{Ass}_{U/S}(\varphi^*\mathcal{F}) \cap U_s$.
Ainsi $x = \varphi(u) \in \text{Ass}_{X/S}(\mathcal{F})$
est au-dessus de $s$, ce qui démontre le lemme.
\end{proof}

\begin{lemma}
\label{lemma-contains-relative-ass-after-base-change}
Soit $Y$ un espace algébrique sur un schéma $S$. Soit $g : X' \to X$ un
morphisme d'espaces algébriques sur $Y$, où $X$ est localement de type fini sur $Y$.
Soit $\mathcal{F}$ un $\mathcal{O}_X$-module quasi-cohérent.
Si $\text{Ass}_{X/Y}(\mathcal{F}) \subset g(|X'|)$, alors, pour tout morphisme
$Z \to Y$, on a $\text{Ass}_{X_Z/Z}(\mathcal{F}_Z) \subset g_Z(|X'_Z|)$.
\end{lemma}

\begin{proof}
D'après Propriétés des espaces, Lemme \ref{spaces-properties-lemma-points-cartesian},
l'application $|X'_Z| \to |X_Z| \times_{|X|} |X'|$ est surjective puisque
$X'_Z$ est égal à $X_Z \times_X X'$.
D'après Diviseurs sur les espaces, Lemme
\ref{spaces-divisors-lemma-base-change-relative-assassin}
l'application $|X_Z| \to |X|$ envoie $\text{Ass}_{X_Z/Z}(\mathcal{F}_Z)$
dans $\text{Ass}_{X/Y}(\mathcal{F})$. Le lemme en résulte.
\end{proof}

\begin{lemma}
\label{lemma-pure-on-top}
Soit $Y$ un espace algébrique sur un schéma $S$. Soit $g : X' \to X$ un
morphisme \'etale d'espaces algébriques sur $Y$. Supposons que les morphismes
structuraux $X' \to Y$ et $X \to Y$ soient décents et de type fini.
Soit $\mathcal{F}$ un $\mathcal{O}_X$-module quasi-cohérent de type fini.
Soit $y \in |Y|$. Posons $F = f^{-1}(\{y\}) \subset |X|$.
\begin{enumerate}
\item Si $\text{Ass}_{X/Y}(\mathcal{F}) \subset g(|X'|)$
et si $g^*\mathcal{F}$ est (universellement) pur au-dessus de $y$, alors
$\mathcal{F}$ est (universellement) pur au-dessus de $y$.
\item Si $\mathcal{F}$ est pur au-dessus de $y$, si $g(|X'|)$ contient $F$, et si
$Y$ est affine local de point fermé $y$, alors
$\text{Ass}_{X/Y}(\mathcal{F}) \subset g(|X'|)$.
\item Si $\mathcal{F}$ est pur au-dessus de $y$, si $\mathcal{F}$ est plat
en tout point de $F$, si $g(|X'|)$ contient
$\text{Ass}_{X/Y}(\mathcal{F}) \cap F$, et si $Y$ est affine local
de point fermé $y$, alors
$\text{Ass}_{X/Y}(\mathcal{F}) \subset g(|X'|)$.
\item Ajouter d'autres assertions ici.
\end{enumerate}
\end{lemma}

\begin{proof}
Les hypothèses sur $X \to Y$ et $X' \to Y$ garantissent que
nous pouvons appliquer les résultats des Sections \ref{section-impure} et
\ref{section-pure}
à ces morphismes et aux faisceaux $\mathcal{F}$ et $g^*\mathcal{F}$.
Comme $g$ est \'etale, nous voyons que
$\text{Ass}_{X'/Y}(g^*\mathcal{F})$
est l'image réciproque de $\text{Ass}_{X/Y}(\mathcal{F})$,
et cela reste vrai après tout changement de base.

\medskip\noindent
Preuve de (1). Supposons $\text{Ass}_{X/Y}(\mathcal{F}) \subset g(|X'|)$.
Supposons que $(T \to Y, t' \leadsto t, \xi)$
soit une impureté de $\mathcal{F}$ au-dessus de $y$. Comme
$\text{Ass}_{X_T/T}(\mathcal{F}_T) \subset g_T(|X'_T|)$ d'après
le Lemme \ref{lemma-contains-relative-ass-after-base-change}, nous pouvons
choisir
un point $\xi' \in |X'_T|$ s'envoyant sur $\xi$. D'après ce qui précède,
$(T \to Y, t' \leadsto t, \xi')$ est une impureté de
$g^*\mathcal{F}$ au-dessus de $y'$. Cela prouve (1).

\medskip\noindent
Preuve de (2). Cela résulte de ce que $g(|X'|)$ est ouvert
dans $|X|$ et de ce que, par pureté, tout point de
$\text{Ass}_{X/Y}(\mathcal{F})$ se spécialise en un point de $F$.

\medskip\noindent
Preuve de (3). Cela résulte de ce que $g(|X'|)$ est ouvert
dans $|X|$ et de ce que la pureté, combinée au
Lemme \ref{lemma-associated-point-specializes}, assure que tout point de
$\text{Ass}_{X/Y}(\mathcal{F})$ se spécialise en un point de
$\text{Ass}_{X/Y}(\mathcal{F}) \cap F$.
\end{proof}

\begin{lemma}
\label{lemma-finite-type-flat-pure-along-fibre-is-universal}
Soit $S$ un schéma. Soit $f : X \to Y$ un morphisme d'espaces
algébriques sur $S$.
Soit $\mathcal{F}$ un $\mathcal{O}_X$-module quasi-cohérent.
Soit $y \in |Y|$.
Supposons que
\begin{enumerate}
\item $f$ soit décent et de type fini,
\item $\mathcal{F}$ soit de type fini,
\item $\mathcal{F}$ soit plat sur $Y$ en tout point au-dessus de $y$, et
\item $\mathcal{F}$ soit pur au-dessus de $y$.
\end{enumerate}
Alors $\mathcal{F}$ est universellement pur au-dessus de $y$.
\end{lemma}

\begin{proof}
Considérons le morphisme $\Spec(\mathcal{O}_{Y, \overline{y}}) \to Y$.
C'est un morphisme plat partant du spectre d'un anneau local
strictement hensélien et envoyant le point fermé sur $y$.
Le Lemme \ref{lemma-flat-descend-pure} nous ramène au cas
décrit dans le paragraphe suivant.

\medskip\noindent
Supposons que $Y$ soit le spectre d'un anneau local strictement hensélien $R$
de point fermé $y$.
D'après le Lemme \ref{lemma-finite-type-flat-along-fibre-free-variant},
il existe un morphisme \'etale $g : X' \to X$ tel que
$g(|X'|) \supset |X_y|$, avec $X'$ affine et
$\Gamma(X', g^*\mathcal{F})$ est un module libre sur $R$.
Alors $g^*\mathcal{F}$ est universellement pur relativement à $Y$, voir
Compléments sur la platitude, Lemme \ref{flat-lemma-affine-locally-projective-pure}.
Il suffit donc de démontrer que
$g(|X'|)$ contient $\text{Ass}_{X/Y}(\mathcal{F})$,
d'après le Lemme \ref{lemma-pure-on-top}, partie (1).
Cela résulte à son tour du
Lemme \ref{lemma-pure-on-top}, partie (2).
\end{proof}

\begin{lemma}
\label{lemma-finite-type-flat-pure-is-universal}
Soit $S$ un schéma.
Soit $f : X \to Y$ un morphisme décent et de type fini d'espaces
algébriques sur $S$.
Soit $\mathcal{F}$ un $\mathcal{O}_X$-module quasi-cohérent de type fini.
Supposons $\mathcal{F}$ plat sur $Y$. Dans ce cas,
$\mathcal{F}$ est pur relativement à $Y$ si et seulement si $\mathcal{F}$
est universellement pur relativement à $Y$.
\end{lemma}

\begin{proof}
C'est une conséquence immédiate du
Lemme \ref{lemma-finite-type-flat-pure-along-fibre-is-universal}
et des définitions.
\end{proof}

\begin{lemma}
\label{lemma-universally-separating}
Soit $Y$ un espace algébrique sur un schéma $S$.
Soit $g : X' \to X$ un morphisme plat d'espaces algébriques sur $Y$,
où $X$ est localement de type fini sur $Y$.
Soit $\mathcal{F}$ un $\mathcal{O}_X$-module quasi-cohérent de type fini
qui est plat sur $Y$. Si $\text{Ass}_{X/Y}(\mathcal{F}) \subset g(|X'|)$,
alors l'application canonique
$$
\mathcal{F} \longrightarrow g_*g^*\mathcal{F}
$$
est injective, et le reste après tout changement de base.
\end{lemma}

\begin{proof}
La dernière assertion signifie que $\mathcal{F}_Z \to (g_Z)_*g_Z^*\mathcal{F}_Z$
est injective pour tout morphisme $Z \to Y$. Comme l'hypothèse sur
l'assassin relatif est préservée par changement de base
(Lemme \ref{lemma-contains-relative-ass-after-base-change}),
il suffit de démontrer l'injectivité de la flèche affichée.

\medskip\noindent
Soit $\mathcal{K} = \Ker(\mathcal{F} \to g_*g^*\mathcal{F})$.
Notre but est de démontrer que $\mathcal{K} = 0$.
Pour cela, il suffit de démontrer que
$\text{WeakAss}_X(\mathcal{K}) = \emptyset$, voir
Diviseurs sur les espaces, Lemme \ref{spaces-divisors-lemma-weakly-ass-zero}.
Nous avons
$\text{WeakAss}_X(\mathcal{K}) \subset \text{WeakAss}_X(\mathcal{F})$, voir
Diviseurs sur les espaces, Lemme \ref{spaces-divisors-lemma-ses-weakly-ass}.
Comme $\mathcal{F}$ est plat, le
Lemme \ref{lemma-bourbaki-finite-type-general-base}
donne $\text{WeakAss}_X(\mathcal{F}) \subset \text{Ass}_{X/Y}(\mathcal{F})$.
Par hypothèse, tout point $x$ de $\text{Ass}_{X/Y}(\mathcal{F})$
est l'image d'un certain $x' \in |X'|$. Puisque $g$ est plat, l'homomorphisme
d'anneaux locaux
$\mathcal{O}_{X, \overline{x}} \to \mathcal{O}_{X', \overline{x}'}$
est fidèlement plat ; par conséquent, l'application
$$
\mathcal{F}_{\overline{x}}
\longrightarrow
(g^*\mathcal{F})_{\overline{x}'} =
\mathcal{F}_{\overline{x}} \otimes_{\mathcal{O}_{X, \overline{x}}}
\mathcal{O}_{X', \overline{x}'}
$$
est injective (voir
Algèbre, Lemme \ref{algebra-lemma-faithfully-flat-universally-injective}).
Puisque la flèche affichée se factorise par
$\mathcal{F}_{\overline{x}} \to (g_*g^*\mathcal{F})_{\overline{x}}$,
nous en concluons que
$\mathcal{K}_{\overline{x}} = 0$. Ainsi $x$ ne peut être un point faiblement
associé de $\mathcal{K}$, ce qui donne le résultat.
\end{proof}









\section{Foncteurs d'aplatissement}
\label{section-F-zero}

\noindent
Cette section est l'analogue de
Compléments sur la platitude, Section \ref{flat-section-flattening-functors}.
Nous conseillons au lecteur de l'omettre en première lecture.

\begin{situation}
\label{situation-iso}
Soit $S$ un schéma.
Soit $f : X \to B$ un morphisme d'espaces algébriques sur $S$.
Soit $u : \mathcal{F} \to \mathcal{G}$ un homomorphisme de
$\mathcal{O}_X$-modules quasi-cohérents. Pour tout schéma $T$ sur
$B$, nous noterons $u_T : \mathcal{F}_T \to \mathcal{G}_T$ le
changement de base de $u$ à $T$ ; autrement dit, $u_T$ est l'image réciproque
de $u$ par le morphisme de projection $X_T = X \times_B T \to X$.
Dans cette situation, nous pouvons considérer le foncteur
\begin{equation}
\label{equation-iso}
F_{iso} : (\Sch/B)^{opp} \longrightarrow \textit{Sets}, \quad
T \longrightarrow \left\{
\begin{matrix}
\{*\} & \text{si }u_T\text{ est un isomorphisme}, \\
\emptyset & \text{sinon.}
\end{matrix}
\right.
\end{equation}
Il existe des variantes $F_{inj}$, $F_{surj}$, $F_{zero}$ pour lesquelles on exige que
$u_T$ soit injectif, surjectif ou nul.
\end{situation}

\noindent
Dans la Situation \ref{situation-iso}, nous considérons parfois les foncteurs
$F_{iso}$, $F_{inj}$, $F_{surj}$ et $F_{zero}$ comme des foncteurs
$(\Sch/S)^{opp} \to \textit{Sets}$ munis d'un morphisme
$F_{iso} \to B$, $F_{inj} \to B$, $F_{surj} \to B$ et $F_{zero} \to B$.
En effet, si $T$ est un schéma sur $S$, un élément $h \in F_{iso}(T)$
est un morphisme $h : T \to B$
tel que le changement de base de $u$ par $h$ soit un isomorphisme.
En particulier, lorsque nous disons
que $F_{iso}$ est un espace algébrique, nous entendons que le foncteur
correspondant $(\Sch/S)^{opp} \to \textit{Sets}$ est un espace algébrique.

\begin{lemma}
\label{lemma-iso-sheaf}
Dans la Situation \ref{situation-iso}.
Chacun des foncteurs $F_{iso}$, $F_{inj}$, $F_{surj}$, $F_{zero}$
vérifie la propriété de faisceau pour la topologie fpqc.
\end{lemma}

\begin{proof}
Soit $\{T_i \to T\}_{i \in I}$ un recouvrement fpqc de schémas sur $B$.
Posons $X_i = X_{T_i} = X \times_S T_i$ et $u_i = u_{T_i}$.
Remarquons que $\{X_i \to X_T\}_{i \in I}$ est un recouvrement fpqc de $X_T$, voir
Topologies sur les espaces, Lemme \ref{spaces-topologies-lemma-fpqc}.
En particulier, pour tout $x \in |X_T|$, il existe $i \in I$
et $x_i \in |X_i|$ s'envoyant sur $x$. Comme
$\mathcal{O}_{X_T, \overline{x}} \to \mathcal{O}_{X_i, \overline{x_i}}$
est plat, donc fidèlement plat (voir
Morphismes d'espaces, Section \ref{spaces-morphisms-section-flat}),
nous concluons que $(u_i)_{x_i}$ est injectif, surjectif, bijectif ou nul
si et seulement si $(u_T)_x$ est injectif, surjectif, bijectif ou nul.
Le lemme en résulte.
\end{proof}

\begin{lemma}
\label{lemma-iso-go-up}
Dans la Situation \ref{situation-iso}, soit $X' \to X$ un morphisme plat
d'espaces algébriques. Notons $u' : \mathcal{F}' \to \mathcal{G}'$
l'image réciproque de $u$ sur $X'$. Notons $F'_{iso}$, $F'_{inj}$, $F'_{surj}$,
$F'_{zero}$ les foncteurs sur $\Sch/B$ associés à $u'$.
\begin{enumerate}
\item Si $\mathcal{G}$ est de type fini et si l'image de $|X'| \to |X|$
contient le support de $\mathcal{G}$, alors $F_{surj} = F'_{surj}$
et $F_{zero} = F'_{zero}$.
\item Si $\mathcal{F}$ est de type fini et si l'image de $|X'| \to |X|$
contient le support de $\mathcal{F}$, alors $F_{inj} = F'_{inj}$
et $F_{zero} = F'_{zero}$.
\item Si $\mathcal{F}$ et $\mathcal{G}$ sont de type fini et si l'image de
$|X'| \to |X|$ contient les supports de $\mathcal{F}$ et de $\mathcal{G}$,
alors $F_{iso} = F'_{iso}$.
\end{enumerate}
\end{lemma}

\begin{proof}
Soit $v : \mathcal{H} \to \mathcal{E}$ une application de modules
quasi-cohérents sur un espace algébrique $Y$, et soit $\varphi : Y' \to Y$ un
morphisme plat surjectif d'espaces algébriques. Alors $v$ est
un isomorphisme, injectif, surjectif ou nul si et seulement si $\varphi^*v$ l'est.
En effet,
pour tout $y \in |Y|$, il existe $y' \in |Y'|$, et l'homomorphisme
d'anneaux locaux
$\mathcal{O}_{Y, \overline{y}} \to \mathcal{O}_{Y', \overline{y'}}$
est fidèlement plat (voir
Morphismes d'espaces, Section \ref{spaces-morphisms-section-flat}).
Bien entendu, pour vérifier l'injectivité ou la nullité, il suffit de considérer
les points du support de $\mathcal{H}$ ; pour vérifier la
surjectivité, il suffit de considérer les points du support de $\mathcal{E}$.
De plus, sous les hypothèses de type fini de l'énoncé du
lemme, la formation des supports commute au changement de base, voir
Morphismes d'espaces, Lemme \ref{spaces-morphisms-lemma-support-finite-type}.
Le lemme est donc clair.
\end{proof}

\noindent
Rappelons que nous avons défini le support schématique d'un module
quasi-cohérent de type fini dans Morphismes d'espaces, Définition
\ref{spaces-morphisms-definition-scheme-theoretic-support}.

\begin{lemma}
\label{lemma-iso-limits}
Dans la Situation \ref{situation-iso}.
\begin{enumerate}
\item Si $\mathcal{G}$ est de type fini et si le support schématique
de $\mathcal{G}$ est quasi-compact sur $B$, alors $F_{surj}$ commute aux
limites.
\item Si $\mathcal{F}$ est de type fini et si le support schématique
de $\mathcal{F}$ est quasi-compact sur $B$, alors
$F_{zero}$ commute aux limites.
\item Si $\mathcal{F}$ est de type fini,
si $\mathcal{G}$ est de présentation finie et si les
supports schématiques de $\mathcal{F}$ et de $\mathcal{G}$ sont
quasi-compacts sur $B$, alors $F_{iso}$ commute aux limites.
\end{enumerate}
\end{lemma}

\begin{proof}
Preuve de (1). Soit $i : Z \to X$ le support schématique de
$\mathcal{G}$, et considérons $\mathcal{G}$ comme un module quasi-cohérent
de type fini sur $Z$. Nous pouvons remplacer $X$ par $Z$ et $u$ par l'application
$i^*\mathcal{F} \to \mathcal{G}$ (détails omis). Nous pouvons donc supposer que
$f$ est quasi-compact et $\mathcal{G}$ de type fini.
Soit $T = \lim_{i \in I} T_i$ une limite dirigée de $B$-schémas affines,
et supposons $u_T$ surjectif.
Posons $X_i = X_{T_i} = X \times_S T_i$ et
$u_i = u_{T_i} : \mathcal{F}_i = \mathcal{F}_{T_i}
\to \mathcal{G}_i = \mathcal{G}_{T_i}$.
Pour démontrer (1), il faut montrer que $u_i$ est surjectif pour un certain $i$.
Choisissons $0 \in I$ et remplaçons $I$ par $\{i \mid i \geq 0\}$.
Puisque $f$ est quasi-compact, $X_0$ est quasi-compact.
Nous pouvons donc choisir un morphisme \'etale surjectif $\varphi_0 : W_0 \to X_0$,
où $W_0$ est un schéma affine. Posons $W = W_0 \times_{T_0} T$
et $W_i = W_0 \times_{T_0} T_i$ pour $i \geq 0$. Ce
sont des schémas affines munis
de morphismes \'etales surjectifs $\varphi : W \to X_T$ et
$\varphi_i : W_i \to X_i$. Remarquons que $W = \lim W_i$.
Ainsi $\varphi^*u_T$ est surjectif, et il suffit de démontrer que
$\varphi_i^*u_i$ est surjectif pour un certain $i$. Nous avons donc ramené
le problème au cas affine, qui est
Algèbre, Lemme \ref{algebra-lemma-module-map-property-in-colimit}, partie (2).

\medskip\noindent
Preuve de (2). Supposons que $\mathcal{F}$ soit de type fini, de support
schématique $Z \subset B$ quasi-compact sur $B$. Soit $T = \lim_{i \in I} T_i$
une limite dirigée de $B$-schémas affines, et supposons $u_T$ nul.
Posons $X_i = T_i \times_B X$ et notons $u_i : \mathcal{F}_i \to \mathcal{G}_i$
l'image réciproque. Choisissons $0 \in I$ et remplaçons $I$ par
$\{i \mid i \geq 0\}$. Posons $Z_0 = Z \times_X X_0$. D'après
Morphismes d'espaces, Lemme \ref{spaces-morphisms-lemma-support-finite-type},
le support de $\mathcal{F}_i$ est $|Z_0|$. Puisque $|Z_0|$ est quasi-compact,
on peut trouver un schéma affine $W_0$ et un morphisme \'etale $W_0 \to X_0$
tels que $|Z_0| \subset \Im(|W_0| \to |X_0|)$.
Posons $W = W_0 \times_{T_0} T$ et $W_i = W_0 \times_{T_0} T_i$ pour $i \geq 0$.
Ce sont des schémas affines munis
de morphismes \'etales $\varphi : W \to X_T$ et
$\varphi_i : W_i \to X_i$. Remarquons que $W = \lim W_i$
et que les supports de $\mathcal{F}_T$ et de $\mathcal{F}_i$
sont contenus respectivement dans les images de $|W| \to |X_T|$ et $|W_i| \to |X_i|$.
Alors $\varphi^*u_T$ est injectif, et il suffit de démontrer que
$\varphi_i^*u_i$ est injectif pour un certain $i$.
Nous avons ainsi ramené le problème au cas affine, qui est
Algèbre, Lemme \ref{algebra-lemma-module-map-property-in-colimit}, partie (1).

\medskip\noindent
Preuve de (3). On peut la démontrer exactement comme dans les
deux paragraphes précédents, en utilisant
Algèbre, Lemme \ref{algebra-lemma-module-map-property-in-colimit}, partie (3).
On peut aussi la déduire de (1) et (2) comme suit.
Soit $T = \lim_{i \in I} T_i$ une limite dirigée de $B$-schémas affines,
et supposons que $u_T$ soit un isomorphisme. D'après (1), il existe
un indice $0 \in I$ tel que $u_{T_0}$ soit surjectif. Posons
$\mathcal{K} = \Ker(u_{T_0})$ et considérons l'application de modules quasi-cohérents
$v : \mathcal{K} \to \mathcal{F}_{T_0}$. Pour $i \geq 0$, le changement de
base $v_{T_i}$ est nul si et seulement si $u_i$ est un isomorphisme. De plus,
$v_T$ est nul. Puisque $\mathcal{G}_{T_0}$
est de présentation finie, que $\mathcal{F}_{T_0}$ est de type fini et que
$u_{T_0}$ est surjectif, nous concluons que $\mathcal{K}$ est de type fini
(Modules sur les sites, Lemme
\ref{sites-modules-lemma-kernel-surjection-finite-onto-finite-presentation}).
Il est clair que le support de $\mathcal{K}$ est contenu dans celui
de $\mathcal{F}_{T_0}$, qui est quasi-compact sur $T_0$.
On peut donc appliquer (2) pour voir que $v_{T_i}$ est nul pour un certain $i$.
\end{proof}

\begin{lemma}
\label{lemma-relate-zero-iso}
Dans la Situation \ref{situation-iso}, supposons donnée une suite exacte
$$
\mathcal{F} \xrightarrow{u} \mathcal{G} \xrightarrow{v} \mathcal{H} \to 0
$$
Alors $F_{v, iso} = F_{u, zero}$, avec les notations évidentes.
\end{lemma}

\begin{proof}
Puisque l'image réciproque est exacte à droite, nous voyons que
$\mathcal{F}_T \to \mathcal{G}_T \to \mathcal{H}_T \to 0$
est exacte pour tout schéma $T$ sur $B$. Ainsi $u_T$ est
surjectif si et seulement si $v_T$ est un isomorphisme.
\end{proof}

\begin{lemma}
\label{lemma-relate-zero-affine}
Dans la Situation \ref{situation-iso}, supposons donnés un morphisme affine
$i : Z \to X$ et un $\mathcal{O}_Z$-module quasi-cohérent $\mathcal{H}$
tels que $\mathcal{G} = i_*\mathcal{H}$. Soit
$v : i^*\mathcal{F} \to \mathcal{H}$ l'application adjointe de $u$.
Alors
\begin{enumerate}
\item $F_{v, zero} = F_{u, zero}$, et
\item si $i$ est une immersion fermée, alors $F_{v, surj} = F_{u, surj}$.
\end{enumerate}
\end{lemma}

\begin{proof}
Soit $T$ un schéma sur $B$. Notons $i_T : Z_T \to X_T$
le changement de base de $i$, et $\mathcal{H}_T$ l'image réciproque de $\mathcal{H}$
sur $Z_T$. Observons que $(i^*\mathcal{F})_T = i_T^*\mathcal{F}_T$
et $i_{T, *}\mathcal{H}_T = (i_*\mathcal{H})_T$.
La première assertion résulte de la commutativité des images réciproques,
et la seconde de Cohomologie des espaces, Lemme
\ref{spaces-cohomology-lemma-affine-base-change}.
Nous voyons donc que $u_T$ et $v_T$ sont eux aussi des applications adjointes.
Ainsi $u_T = 0$ si et seulement si $v_T = 0$. Cela démontre (1).
Dans le cas (2), nous voyons que $u_T$ est surjectif si et seulement si
$v_T$ est surjectif, car $u_T$ se factorise sous la forme
$$
\mathcal{F}_T \to
i_{T, *}i_T^*\mathcal{F}_T \xrightarrow{i_{T, *}v_T} i_{T, *}\mathcal{H}_T
$$
et parce que $i_{T, *}$ est un foncteur exact
qui plonge pleinement fidèlement la catégorie des modules quasi-cohérents sur
$Z_T$ dans la catégorie des $\mathcal{O}_{X_T}$-modules quasi-cohérents.
Voir Morphismes d'espaces, Lemme \ref{spaces-morphisms-lemma-i-star-equivalence}.
\end{proof}

\begin{lemma}
\label{lemma-relate-zero-affine-push}
Dans la Situation \ref{situation-iso}, supposons donné un morphisme affine
$g : X \to X'$. Posons $u' = f_*u : f_*\mathcal{F} \to f_*\mathcal{G}$.
Alors $F_{u, iso} = F_{u', iso}$, $F_{u, inj} = F_{u', inj}$,
$F_{u, surj} = F_{u', surj}$ et $F_{u, zero} = F_{u', zero}$.
\end{lemma}

\begin{proof}
D'après Cohomologie des espaces, Lemme
\ref{spaces-cohomology-lemma-affine-base-change}
nous avons $g_{T, *}u_T = u'_T$.
De plus, $g_{T, *} : \QCoh(\mathcal{O}_{X_T}) \to \QCoh(\mathcal{O}_X)$
est un foncteur fidèle et exact qui reflète les isomorphismes, les applications
injectives et les applications surjectives.
\end{proof}

\begin{situation}
\label{situation-flat}
Soit $S$ un schéma.
Soit $f : X \to Y$ un morphisme d'espaces algébriques sur $S$.
Soit $\mathcal{F}$ un $\mathcal{O}_X$-module quasi-cohérent.
Pour tout schéma $T$ sur $Y$, nous noterons $\mathcal{F}_T$ le
changement de base de $\mathcal{F}$ à $T$ ; autrement dit, $\mathcal{F}_T$
est l'image réciproque de $\mathcal{F}$ par le morphisme de projection
$X_T = X \times_Y T \to X$. Puisque tout changement de base d'un module plat
est plat, nous obtenons un foncteur
\begin{equation}
\label{equation-flat}
F_{flat} : (\Sch/Y)^{opp} \longrightarrow \textit{Sets}, \quad
T \longrightarrow \left\{
\begin{matrix}
\{*\} & \text{si } \mathcal{F}_T \text{ est plat sur }T, \\
\emptyset & \text{sinon.}
\end{matrix}
\right.
\end{equation}
\end{situation}

\noindent
Dans la Situation \ref{situation-flat}, nous considérons parfois $F_{flat}$ comme un foncteur
$(\Sch/S)^{opp} \to \textit{Sets}$ muni d'un morphisme
$F_{flat} \to Y$. En effet, si $T$ est un schéma sur $S$,
un élément $h \in F_{flat}(T)$ est un morphisme $h : T \to Y$
tel que le changement de base de $\mathcal{F}$ par $h$ soit plat sur $T$.
En particulier, lorsque nous disons
que $F_{flat}$ est un espace algébrique, nous entendons que le foncteur
correspondant $(\Sch/S)^{opp} \to \textit{Sets}$ est un espace algébrique.

\begin{lemma}
\label{lemma-flat}
Dans la Situation \ref{situation-flat}.
\begin{enumerate}
\item Le foncteur $F_{flat}$ vérifie la propriété de faisceau pour la topologie fpqc.
\item Si $f$ est quasi-compact et localement de présentation finie
et si $\mathcal{F}$ est de présentation finie, alors le foncteur
$F_{flat}$ commute aux limites.
\end{enumerate}
\end{lemma}

\begin{proof}
La partie (1) résulte de l'assertion suivante : si $T' \to T$ est un morphisme
plat surjectif d'espaces algébriques sur $Y$, alors
$\mathcal{F}_{T'}$ est plat sur $T'$ si et seulement si
$\mathcal{F}_T$ est plat sur $T$, voir
Morphismes d'espaces, Lemme \ref{spaces-morphisms-lemma-base-change-module-flat}.
La partie (2) résulte de
Limites d'espaces, Lemme \ref{spaces-limits-lemma-descend-flat},
si $f$ est en outre quasi-séparé (c'est-à-dire si $f$ est de présentation finie).
Dans le cas général, on se ramène d'abord au cas où la
base est affine, puis on recouvre $X$ par un nombre fini d'ouverts affines
pour se ramener au cas quasi-séparé. Détails omis.
\end{proof}










\section{Annulation d'une application}
\label{section-zero-map}

\noindent
Cette section n'a pas d'analogue dans le chapitre correspondant sur les schémas.

\begin{situation}
\label{situation-somewhat-closed}
Soit $S = \Spec(R)$ un schéma affine. Soit $X$ un espace algébrique sur
$S$. Soit $u : \mathcal{F} \to \mathcal{G}$ une application de
$\mathcal{O}_X$-modules quasi-cohérents. Supposons $\mathcal{G}$ plat sur $S$.
\end{situation}

\begin{lemma}
\label{lemma-F-zero-somewhat-closed}
Dans la Situation \ref{situation-somewhat-closed}. Soit $T \to S$
un morphisme quasi-compact de schémas tel que le changement de base $u_T$ soit
nul. Alors il existe un sous-schéma fermé $Z \subset S$ tel que
(a) $T \to S$ se factorise par $Z$ et (b) le changement de base $u_Z$ soit nul.
Si $\mathcal{F}$ est un $\mathcal{O}_X$-module de type fini et si
le support schématique de $\mathcal{F}$ est quasi-compact,
on peut prendre $Z \to S$ de présentation finie.
\end{lemma}

\begin{proof}
Soit $U \to X$ un morphisme \'etale surjectif d'espaces algébriques,
où $U = \coprod U_i$ est une réunion disjointe de schémas affines (voir
Propriétés des espaces, Lemme
\ref{spaces-properties-lemma-cover-by-union-affines}).
D'après le Lemme \ref{lemma-iso-go-up}, nous pouvons
remplacer $X$ par $U$. Autrement dit, nous pouvons supposer que $X = \coprod X_i$
est une réunion disjointe de schémas affines $X_i$. Supposons que nous sachions démontrer
le lemme pour $u_i = u|_{X_i}$. Nous obtenons alors un sous-schéma fermé
$Z_i \subset S$ tel que $T \to S$ se factorise par $Z_i$ et que
$u_{i, Z_i}$ soit nul. Si
$Z_i = \Spec(R/I_i) \subset \Spec(R) = S$, alors prendre
$Z = \Spec(R/\sum I_i)$ convient. Nous pouvons donc supposer que
$X = \Spec(A)$ est affine.

\medskip\noindent
Choisissons un recouvrement ouvert affine fini $T = T_1 \cup \ldots \cup T_m$.
Il est clair que nous pouvons remplacer $T$ par $\coprod_{j = 1, \ldots, m} T_j$.
Nous pouvons donc supposer $T$ affine. Écrivons $T = \Spec(R')$.
Soit $u : M \to N$ l'homomorphisme de $A$-modules
correspondant à $u : \mathcal{F} \to \mathcal{G}$.
Alors $N$ est un $R$-module plat puisque $\mathcal{G}$ est plat sur $S$.
L'hypothèse du lemme signifie que le composé
$$
M \otimes_R R' \to N \otimes_R R'
$$
est nul. Soit $z \in M$. D'après le théorème de Lazard
(Algèbre, Théorème \ref{algebra-theorem-lazard}) et puisque
$\otimes$ commute aux colimites, on peut trouver un $R$-module libre
$F_z$, un élément $\tilde z \in F_z$ et une application $F_z \to N$ tels que
$u(z)$ soit l'image de $\tilde z$ et que $\tilde z$ s'envoie sur zéro dans
$F_z \otimes_R R'$. Choisissons une base $\{e_{z, \alpha}\}$ de $F_z$ et écrivons
$\tilde z = \sum f_{z, \alpha} e_{z, \alpha}$ avec $f_{z, \alpha} \in R$.
Soit $I \subset R$ l'idéal engendré par les éléments $f_{z, \alpha}$,
où $z$ parcourt tous les éléments de $M$.
Par construction, $I$ s'envoie sur zéro dans $R'$ et les éléments $\tilde z$
s'envoient sur zéro dans $F_z/IF_z$, donc dans $N/IN$. Ainsi $Z = \Spec(R/I)$
résout le problème dans ce cas.

\medskip\noindent
Supposons que $\mathcal{F}$ soit de type fini, de support schématique
quasi-compact. Écrivons $Z = \Spec(R/I)$.
Écrivons $I = \bigcup I_\lambda$ comme réunion filtrante d'idéaux de type fini.
Posons $Z_\lambda = \Spec(R/I_\lambda)$, de sorte que $Z = \colim Z_\lambda$.
Puisque $u_Z$ est nul, nous voyons que $u_{Z_\lambda}$ est nul
pour un certain $\lambda$, d'après le Lemme \ref{lemma-iso-limits}.
Cela achève la démonstration du lemme.
\end{proof}

\begin{lemma}
\label{lemma-F-zero-module-map}
Soit $A$ un anneau. Soit $u : M \to N$ une application de $A$-modules.
Si $N$ est projectif comme $A$-module, alors il existe un idéal
$I \subset A$ tel que, pour tout homomorphisme d'anneaux $\varphi : A \to B$,
les conditions suivantes soient équivalentes :
\begin{enumerate}
\item $u \otimes 1 : M \otimes_A B \to N \otimes_A B$ est nul, et
\item $\varphi(I) = 0$.
\end{enumerate}
\end{lemma}

\begin{proof}
Puisque $N$ est projectif, on peut trouver un $A$-module projectif $C$
tel que $F = N \oplus C$ soit un $R$-module libre.
En remplaçant $u$ par $u \oplus 1 : F = M \oplus C \to N \oplus C$,
nous voyons que nous pouvons supposer $N$ libre. Dans ce cas, soit $I$
l'idéal de $A$ engendré par les coefficients de tous les éléments de
$\Im(u)$ relativement à une base (fixée) de $N$.
\end{proof}

\begin{lemma}
\label{lemma-F-zero-somewhat-closed-points}
Dans la Situation \ref{situation-somewhat-closed}.
Soit $T \subset S$ une partie. Soit $s \in S$ dans l'adhérence de $T$.
Pour $t \in T$, soit $u_t$ l'image réciproque de $u$ sur $X_t$,
et soit $u_s$ l'image réciproque de $u$ sur $X_s$.
Si $X$ est localement de présentation finie sur $S$,
si $\mathcal{G}$ est de présentation finie\footnote{Il suffirait
que $X$ soit localement de type fini sur $S$
et que $\mathcal{G}$ soit de présentation finie relativement à $S$,
mais cette notion n'a pas encore été définie dans le cadre
des espaces algébriques. La définition pour les schémas est
donnée dans Compléments sur les morphismes, Section
\ref{more-morphisms-section-finite-type-finite-presentation}.}, et si
$u_t = 0$ pour tout $t \in T$, alors $u_s = 0$.
\end{lemma}

\begin{proof}
Vérifier si $u_s$ est nul est une question locale pour la topologie \'etale sur la fibre $X_s$.
Nous pouvons donc choisir un point $x \in |X_s| \subset |X|$ et effectuer la vérification
dans un voisinage \'etale. Choisissons
$$
\xymatrix{
(X, x) \ar[d] & (X', x') \ar[l]^g \ar[d] \\
(S, s) & (S', s') \ar[l]
}
$$
comme dans la Proposition \ref{proposition-finite-presentation-flat-at-point}.
Soit $T' \subset S'$ l'image réciproque de $T$. Observons que
$s'$ appartient à l'adhérence de $T'$, car $S' \to S$ est ouvert.
Nous sommes donc ramenés au problème d'algèbre décrit dans le
paragraphe suivant.

\medskip\noindent
Nous avons une application de $R$-modules $u : M \to N$ telle que $N$ soit projectif
comme $R$-module et telle que
$u_t : M \otimes_R \kappa(t) \to N \otimes_R \kappa(t)$
soit nulle pour tout $t \in T$. Il s'agit de montrer que $u_s = 0$.
Soit $I \subset R$ l'idéal défini dans le Lemme \ref{lemma-F-zero-module-map}.
Alors $I$ s'envoie sur zéro dans $\kappa(t)$ pour tout $t \in T$.
Ainsi $T \subset V(I)$. Puisque $s$ appartient à l'adhérence de $T$,
nous avons $s \in V(I)$. Par conséquent, $u_s = 0$.
\end{proof}

\noindent
Il serait intéressant de trouver une démonstration directe « simple » de l'un des
Lemmes \ref{lemma-F-zero-closed-pure} ou
\ref{lemma-F-zero-closed-proper},
en utilisant des arguments analogues à ceux des
Lemmes \ref{lemma-F-zero-somewhat-closed} et
\ref{lemma-F-zero-somewhat-closed-points}.
Une démonstration « classique » de ce lemme lorsque $f : X \to B$ est un morphisme
projectif et $B$ un schéma noethérien serait la suivante : (a) choisir un faisceau inversible
relativement ample $\mathcal{O}_X(1)$, (b) poser
$u_n : f_*\mathcal{F}(n) \to f_*\mathcal{G}(n)$,
(c) observer que $f_*\mathcal{G}(n)$ est un faisceau localement libre de rang fini
pour tout $n \gg 0$, et (d) que $F_{zero}$ est représenté par le lieu des zéros
de $u_n$ pour un certain $n \gg 0$.

\begin{lemma}
\label{lemma-F-zero-closed-pure}
Dans la Situation \ref{situation-iso}. Supposons que
\begin{enumerate}
\item $f$ soit de présentation finie, et
\item $\mathcal{G}$ soit de présentation finie,
plat sur $B$ et pur relativement à $B$.
\end{enumerate}
Alors $F_{zero}$ est un espace algébrique et $F_{zero} \to B$
est une immersion fermée. Si $\mathcal{F}$ est de type fini, alors
$F_{zero} \to B$ est de présentation finie.
\end{lemma}

\begin{proof}
D'après le Lemme \ref{lemma-finite-type-flat-pure-is-universal},
le module $\mathcal{G}$ est universellement pur relativement à $B$.
Pour démontrer que $F_{zero}$ est un espace algébrique,
il suffit de montrer que $F_{zero} \to B$ est représentable, voir
Espaces, Lemme \ref{spaces-lemma-representable-over-space}.
Soit $B' \to B$ un morphisme, où $B'$ est un schéma, et soit
$u' : \mathcal{F}' \to \mathcal{G}'$ l'image réciproque de $u$ sur $X' = X_{B'}$.
Alors le foncteur associé $F'_{zero}$ est égal à $F_{zero} \times_B B'$.
Cela nous ramène au cas où $B$ est un schéma.

\medskip\noindent
Supposons que $B$ soit un schéma. Nous allons montrer que $F_{zero}$ est représentable
par un sous-schéma fermé de $B$. D'après le Lemme \ref{lemma-iso-sheaf} et
Descente, Lemmes \ref{descent-lemma-closed-immersion} et
\ref{descent-lemma-descent-data-sheaves}
la question est locale pour la topologie \'etale sur $B$. Soit $b \in B$.
Remplaçons d'abord $B$ par un voisinage affine de $b$.
Choisissons un diagramme
$$
\xymatrix{
X \ar[d] & X' \ar[l]^g \ar[d] \\
B & B' \ar[l]
}
$$
et un point $b' \in B'$ s'envoyant sur $b \in B$,
comme dans le Lemme \ref{lemma-finite-presentation-flat-along-fibre}.
Puisque nous travaillons localement pour la topologie \'etale, nous pouvons remplacer
$B$ par $B'$ et supposer que nous disposons d'un diagramme
$$
\xymatrix{
X \ar[rd] & & X' \ar[ll]^g \ar[ld] \\
& B
}
$$
où $B$ et $X'$ sont affines, tel que $\Gamma(X', g^*\mathcal{G})$
soit un $\Gamma(B, \mathcal{O}_B)$-module projectif et que
$g(|X'|) \supset |X_b|$. Soit $U \subset X$ le sous-espace ouvert
tel que $|U| = g(|X'|)$. D'après
Diviseurs sur les espaces, Lemme
\ref{spaces-divisors-lemma-relative-assassin-constructible}, l'ensemble
$$
E = \{t \in B : \text{Ass}_{X_t}(\mathcal{G}_t) \subset |U_t|\} =
\{t \in B : \text{Ass}_{X/B}(\mathcal{G}) \cap |X_t| \subset |U_t|\}
$$
est constructible dans $B$. D'après le Lemme \ref{lemma-pure-on-top}, partie (2),
nous voyons que $E$ contient $\Spec(\mathcal{O}_{B, b})$. D'après
Morphismes, Lemme \ref{morphisms-lemma-constructible-containing-open},
nous voyons que $E$ contient un voisinage ouvert de $b$. Ainsi, après avoir
remplacé $B$ par un voisinage affine plus petit de $b$, nous pouvons supposer que
$\text{Ass}_{X/B}(\mathcal{G}) \subset g(|X'|)$.

\medskip\noindent
Il résulte du Lemme \ref{lemma-universally-separating}
que $u : \mathcal{F} \to \mathcal{G}$ est injectif si et seulement si
$g^*u : g^*\mathcal{F} \to g^*\mathcal{G}$ est injectif, et cela reste
vrai après tout changement de base. Nous nous sommes donc ramenés au cas où,
outre les hypothèses du théorème, $X \to B$ est un morphisme de
schémas affines et $\Gamma(X, \mathcal{G})$ est un
$\Gamma(B, \mathcal{O}_B)$-module projectif. Ce cas résulte immédiatement du
Lemme \ref{lemma-F-zero-module-map}.

\medskip\noindent
Il reste à montrer que $F_{zero} \to B$ est de présentation finie si
$\mathcal{F}$ est de type fini. Cela résulte du
Lemme \ref{lemma-iso-limits}, combiné à
Limites d'espaces, Proposition
\ref{spaces-limits-proposition-characterize-locally-finite-presentation}.
\end{proof}

\begin{lemma}
\label{lemma-F-zero-closed-proper}
Dans la Situation \ref{situation-iso}. Supposons que
\begin{enumerate}
\item $f$ soit localement de présentation finie,
\item $\mathcal{G}$ soit un $\mathcal{O}_X$-module de présentation finie
plat sur $B$,
\item le support de $\mathcal{G}$ soit propre sur $B$.
\end{enumerate}
Alors le foncteur $F_{zero}$ est un espace algébrique et $F_{zero} \to B$
est une immersion fermée. Si $\mathcal{F}$ est de type fini, alors
$F_{zero} \to B$ est de présentation finie.
\end{lemma}

\begin{proof}
Si $f$ est de présentation finie, cela résulte immédiatement des
Lemmes \ref{lemma-F-zero-closed-pure} et \ref{lemma-proper-pure}.
C'est le seul cas intéressant, et nous conseillons au lecteur d'omettre
la suite de la démonstration, qui traite de la possibilité (autorisée
par les hypothèses du présent lemme)
que $f$ ne soit ni quasi-séparé ni quasi-compact.

\medskip\noindent
Soit $i : Z \to X$ le sous-espace fermé défini par l'idéal de Fitting
d'indice zéro de $\mathcal{G}$
(Diviseurs sur les espaces, Section
\ref{spaces-divisors-section-fitting-ideals}).
Alors $Z \to B$ est propre par hypothèse (voir
Catégories dérivées des espaces, Section
\ref{spaces-perfect-section-proper-over-base}).
D'autre part, $i$ est de présentation finie
(Diviseurs sur les espaces, Lemme
\ref{spaces-divisors-lemma-fitting-ideal-of-finitely-presented}, et
Morphismes d'espaces, Lemme
\ref{spaces-morphisms-lemma-closed-immersion-finite-presentation}).
Il existe un $\mathcal{O}_Z$-module quasi-cohérent
$\mathcal{H}$ de type fini tel que $i_*\mathcal{H} = \mathcal{G}$
(Diviseurs sur les espaces, Lemme
\ref{spaces-divisors-lemma-on-subscheme-cut-out-by-Fit-0}).
En fait, $\mathcal{H}$ est de présentation finie comme $\mathcal{O}_Z$-module
d'après Algèbre, Lemme \ref{algebra-lemma-finitely-presented-over-subring}
(détails omis).
Alors $F_{zero}$ est le même que le foncteur $F_{zero}$
pour l'application $i^*\mathcal{F} \to \mathcal{H}$ adjointe de $u$, voir
Lemme \ref{lemma-relate-zero-affine}.
Le faisceau $\mathcal{H}$ est plat sur $B$, car
il en est de même de $\mathcal{G}$ (vérifier sur les germes ; détails omis).
De plus, remarquons que si $\mathcal{F}$ est de type fini,
alors $i^*\mathcal{F}$ est de type fini.
Nous avons donc ramené le lemme au cas
étudié dans le premier paragraphe de la démonstration.
\end{proof}








\section{Aplatissement d'une application}
\label{section-flattening-map}

\noindent
Cette section est l'analogue de
Compléments sur la platitude, Section \ref{flat-section-flattening-map}.
En particulier, le résultat suivant est une variante de
Compléments sur la platitude, Théorème \ref{flat-theorem-flattening-map}.

\begin{theorem}
\label{theorem-flattening-map}
Dans la Situation \ref{situation-iso}, supposons que
\begin{enumerate}
\item $f$ soit de présentation finie,
\item $\mathcal{F}$ soit de présentation finie, plat sur $B$ et
pur relativement à $B$, et
\item $u$ soit surjectif.
\end{enumerate}
Alors $F_{iso}$ est représentable par une immersion fermée $Z \to B$.
De plus, $Z \to S$ est de présentation finie si $\mathcal{G}$ est
de présentation finie.
\end{theorem}

\begin{proof}
Soit $\mathcal{K} = \Ker(u)$ et notons $v : \mathcal{K} \to \mathcal{F}$
l'inclusion. Le Lemme \ref{lemma-relate-zero-iso} donne
$F_{u, iso} = F_{v, zero}$. En appliquant le Lemme \ref{lemma-F-zero-closed-pure}
à $v$, nous voyons que $F_{u, iso} = F_{v, zero}$ est représentable
par un sous-espace fermé de $B$. Remarquons que $\mathcal{K}$ est de type fini
si $\mathcal{G}$ est de présentation finie, voir
Modules sur les sites, Lemme
\ref{sites-modules-lemma-kernel-surjection-finite-onto-finite-presentation}.
Nous obtenons donc également la dernière assertion du théorème.
\end{proof}

\begin{lemma}
\label{lemma-F-iso-closed}
Dans la Situation \ref{situation-iso}. Supposons que
\begin{enumerate}
\item $f$ soit localement de présentation finie,
\item $\mathcal{F}$ soit localement de présentation finie et plat sur $B$,
\item le support de $\mathcal{F}$ soit propre sur $B$, et
\item $u$ soit surjectif.
\end{enumerate}
Alors le foncteur $F_{iso}$ est un espace algébrique et $F_{iso} \to B$
est une immersion fermée. Si $\mathcal{G}$ est de présentation finie, alors
$F_{iso} \to B$ est de présentation finie.
\end{lemma}

\begin{proof}
Soit $\mathcal{K} = \Ker(u)$ et notons $v : \mathcal{K} \to \mathcal{F}$
l'inclusion. Le Lemme \ref{lemma-relate-zero-iso} donne
$F_{u, iso} = F_{v, zero}$. En appliquant le Lemme \ref{lemma-F-zero-closed-proper}
à $v$, nous voyons que $F_{u, iso} = F_{v, zero}$ est représentable
par un sous-espace fermé de $B$. Remarquons que $\mathcal{K}$ est de type fini
si $\mathcal{G}$ est de présentation finie, voir
Modules sur les sites, Lemme
\ref{sites-modules-lemma-kernel-surjection-finite-onto-finite-presentation}.
Nous obtenons donc également la dernière assertion du lemme.
\end{proof}

\noindent
Nous utiliserons le résultat (facile) suivant dans l'étude du foncteur Quot.

\begin{lemma}
\label{lemma-F-surj-open}
Dans la Situation \ref{situation-iso}. Supposons que
\begin{enumerate}
\item $f$ soit localement de présentation finie,
\item $\mathcal{G}$ soit de type fini,
\item le support de $\mathcal{G}$ soit propre sur $B$.
\end{enumerate}
Alors $F_{surj}$ est un espace algébrique et $F_{surj} \to B$
est une immersion ouverte.
\end{lemma}

\begin{proof}
Considérons $\Coker(u)$. Observons que
$\Coker(u_T) = \Coker(u)_T$ pour tout $T/B$.
Remarquons que la formation du support d'un module
quasi-cohérent de type fini commute à l'image réciproque
(Morphismes d'espaces, Lemme \ref{spaces-morphisms-lemma-support-covering}).
Ainsi $F_{surj}$ est représentable par le sous-espace ouvert de $B$
correspondant à l'ouvert
$$
|B| \setminus |f|(\text{Supp}(\Coker(u)))
$$
voir Propriétés des espaces, Lemme \ref{spaces-properties-lemma-open-subspaces}.
C'est un ouvert, car $|f|$ est fermé sur $\text{Supp}(\mathcal{G})$
et $\text{Supp}(\Coker(u))$ est une partie fermée de
$\text{Supp}(\mathcal{G})$.
\end{proof}






\section{Aplatissement dans le cas local}
\label{section-flattening-local}

\noindent
Cette section est l'analogue de Compléments sur la platitude, Section
\ref{flat-section-flattening-local}.

\begin{lemma}
\label{lemma-freebie}
Soit $S$ le spectre d'un anneau local hensélien de point fermé $s$.
Soit $X \to S$ un morphisme d'espaces algébriques
localement de type fini.
Soit $\mathcal{F}$ un $\mathcal{O}_X$-module quasi-cohérent de type fini.
Soit $E \subset |X_s|$ une partie. Il existe un sous-schéma fermé
$Z \subset S$ ayant la propriété suivante : pour tout morphisme de schémas
pointés $(T, t) \to (S, s)$, les conditions suivantes sont équivalentes :
\begin{enumerate}
\item $\mathcal{F}_T$ est plat sur $T$ en tout point de
$|X_t|$ qui s'envoie sur un point de $E \subset |X_s|$, et
\item $\Spec(\mathcal{O}_{T, t}) \to S$ se factorise par $Z$.
\end{enumerate}
De plus, si $X \to S$ est localement de présentation finie,
si $\mathcal{F}$ est de présentation finie et si $E \subset |X_s|$ est
fermé et quasi-compact, alors $Z \to S$ est de présentation finie.
\end{lemma}

\begin{proof}
Choisissons un schéma $U$ et un morphisme \'etale $\varphi : U \to X$.
Soit $E' \subset |U_s|$ l'image réciproque de $E$. Si
$E' \to E$ est surjectif, alors la condition (1) équivaut à ceci :
$(\varphi^*\mathcal{F})_T$ est plat sur $T$ en tout point de
$|U_t|$ qui s'envoie sur un point de $E' \subset |U_t|$.
En choisissant $\varphi$ surjectif, nous sommes ramenés au cas des schémas, à savoir
Compléments sur la platitude, Lemme \ref{flat-lemma-freebie}.
Si $E$ est fermé et quasi-compact, nous pouvons choisir $U$
affine de telle sorte que $E' \to E$ soit surjectif. Alors $E'$ est fermé
et quasi-compact, et la dernière assertion résulte de la
dernière assertion de
Compléments sur la platitude, Lemme \ref{flat-lemma-freebie}.
\end{proof}









\section{Aplatissement universel}
\label{section-flattening-final}

\noindent
Cette section est l'analogue de
Compléments sur la platitude, Section \ref{flat-section-flattening-final}.
Notre but principal est de démontrer le
Lemme \ref{lemma-when-universal-flattening}.
Cependant, nous ne voyons pas comment déduire directement ce résultat
du résultat correspondant pour les schémas.
Nous devons donc reprendre ici une partie du développement.
Commençons toutefois par une définition.

\begin{definition}
\label{definition-flattening}
Soit $S$ un schéma.
Soit $X \to Y$ un morphisme d'espaces algébriques sur $S$.
Soit $\mathcal{F}$ un $\mathcal{O}_X$-module quasi-cohérent.
Nous disons que {\it l'aplatissement universel de $\mathcal{F}$ existe}
si le foncteur $F_{flat}$ défini dans la Situation \ref{situation-flat}
est un espace algébrique.
Nous disons que {\it l'aplatissement universel de $X$ existe}
si l'aplatissement universel de $\mathcal{O}_X$ existe.
\end{definition}

\noindent
Cette définition est quelque peu insatisfaisante, car la définition de
l'aplatissement universel ne coïncide pas ici avec celle utilisée dans le
cas des schémas, puisque nous ignorons si tout monomorphisme d'espaces
algébriques est représentable (Compléments sur les morphismes d'espaces,
Section \ref{spaces-more-morphisms-section-monomorphisms}).
Espérons qu'il n'en résultera jamais de confusion.

\begin{lemma}
\label{lemma-pre-flat-dimension-n}
Soit $S$ un schéma.
Soit $f : X \to Y$ un morphisme d'espaces algébriques
localement de type fini.
Soit $\mathcal{F}$ un $\mathcal{O}_X$-module quasi-cohérent de
type fini. Soit $n \geq 0$. Les conditions suivantes sont équivalentes :
\begin{enumerate}
\item pour un certain diagramme commutatif
$$
\xymatrix{
U \ar[d]_\varphi \ar[r] & V \ar[d] \\
X \ar[r] & Y
}
$$
dont les flèches verticales sont \'etales et surjectives, où $U$ et $V$ sont
des schémas, le faisceau $\varphi^*\mathcal{F}$ est plat sur $V$
en dimensions $\geq n$ (Compléments sur la platitude, Définition
\ref{flat-definition-flat-dimension-n}),
\item pour tout diagramme commutatif
$$
\xymatrix{
U \ar[d]_\varphi \ar[r] & V \ar[d] \\
X \ar[r] & Y
}
$$
dont les flèches verticales sont \'etales, où $U$ et $V$ sont des schémas,
le faisceau $\varphi^*\mathcal{F}$ est plat sur $V$ en dimensions $\geq n$, et
\item pour $x \in |X|$ tel que $\mathcal{F}$ ne soit pas plat en $x$
sur $Y$, le degré de transcendance de $x/f(x)$ est $< n$ (Morphismes d'espaces,
Définition \ref{spaces-morphisms-definition-dimension-fibre}).
\end{enumerate}
Si ces conditions sont satisfaites, elles le restent après tout changement de base $Y' \to Y$.
\end{lemma}

\begin{proof}
Supposons disposer d'un diagramme comme en (1). L'équivalence des
conditions dans Compléments sur la platitude, Lemme \ref{flat-lemma-pre-flat-dimension-n},
montre alors que (1) et (3) sont équivalentes. Mais la condition (3) est héritée
par $\varphi^*\mathcal{F}$ pour tout $U \to V$ comme en (2).
Le résultat pour les schémas montre donc à nouveau que (3) implique (2).
Il implique également l'assertion relative au changement de base.
\end{proof}

\begin{definition}
\label{definition-flat-dimension-n}
Soit $S$ un schéma. Soit $f : X \to Y$ un morphisme d'espaces algébriques
sur $S$ qui est localement de type fini.
Soit $\mathcal{F}$ un $\mathcal{O}_X$-module quasi-cohérent de type fini.
Soit $n \geq 0$.
Nous disons que {\it $\mathcal{F}$ est plat sur $Y$ en dimensions $\geq n$}
si les conditions équivalentes du Lemme \ref{lemma-pre-flat-dimension-n}
sont satisfaites.
\end{definition}

\begin{situation}
\label{situation-flat-dimension-n}
Soit $S$ un schéma. Soit $f : X \to Y$ un morphisme d'espaces algébriques
sur $S$ qui est localement de type fini. Soit $\mathcal{F}$ un
$\mathcal{O}_X$-module quasi-cohérent de type fini. Pour tout schéma $T$ sur $Y$, nous
noterons $\mathcal{F}_T$ le changement de base de $\mathcal{F}$ à $T$ ; autrement dit,
$\mathcal{F}_T$ est l'image réciproque de $\mathcal{F}$ par le morphisme de projection
$X_T = X \times_Y T \to X$. Remarquons que $f_T : X_T \to T$ est de type fini
et que $\mathcal{F}_T$ est un $\mathcal{O}_{X_T}$-module de type fini
(Morphismes d'espaces, Lemme
\ref{spaces-morphisms-lemma-base-change-finite-type}
et
Modules sur les sites, Lemme \ref{sites-modules-lemma-local-pullback}).
Soit $n \geq 0$. La Définition \ref{definition-flat-dimension-n} et le
Lemme \ref{lemma-pre-flat-dimension-n} fournissent un foncteur
\begin{equation}
\label{equation-flat-dimension-n}
F_n : (\Sch/Y)^{opp} \longrightarrow \textit{Sets}, \quad
T \longrightarrow \left\{
\begin{matrix}
\{*\} & \text{si }\mathcal{F}_T\text{ est plat sur }T\text{ en }\dim \geq n, \\
\emptyset & \text{sinon.}
\end{matrix}
\right.
\end{equation}
\end{situation}

\noindent
Dans la Situation \ref{situation-flat-dimension-n}, nous considérons parfois
$F_n$ comme un foncteur $(\Sch/S)^{opp} \to \textit{Sets}$ muni d'un morphisme
$F_n \to Y$. En effet, si $T$ est un schéma sur $S$,
un élément $h \in F_n(T)$ est un morphisme $h : T \to Y$
tel que le changement de base de $\mathcal{F}$ par $h$ soit plat sur $T$
en $\dim \geq n$. En particulier, lorsque nous disons
que $F_n$ est un espace algébrique, nous entendons que le foncteur
correspondant $(\Sch/S)^{opp} \to \textit{Sets}$ est un espace algébrique.

\begin{lemma}
\label{lemma-flat-dimension-n}
Dans la Situation \ref{situation-flat-dimension-n}.
\begin{enumerate}
\item Le foncteur $F_n$ vérifie la propriété de faisceau pour la topologie fpqc.
\item Si $f$ est quasi-compact et localement de présentation finie
et si $\mathcal{F}$ est de présentation finie, alors le foncteur $F_n$
commute aux limites.
\end{enumerate}
\end{lemma}

\begin{proof}
Preuve de (1). Supposons que $\{T_i \to T\}$ soit un recouvrement fpqc
d'un schéma $T$ sur $Y$. Nous devons montrer que, si $F_n(T_i)$ est non vide
pour tout $i$, alors $F_n(T)$ est non vide.
Choisissons un diagramme comme dans la partie (1) du Lemme \ref{lemma-pre-flat-dimension-n}.
Notons $F'_n$ le foncteur correspondant pour
$\varphi^*\mathcal{F}$ et le morphisme $U \to V$.
D'après Compléments sur la platitude, Lemme \ref{flat-lemma-flat-dimension-n},
$F'_n$ vérifie la propriété de faisceau.
Il en est donc de même de $F_n$, car
pour $T \to Y$, on a $F_n(T) = F'_n(V \times_Y T)$
d'après le Lemme \ref{lemma-pre-flat-dimension-n},
et parce que $\{V \times_Y T_i \to V \times_Y T\}$
est un recouvrement fpqc.

\medskip\noindent
Preuve de (2). Supposons que $T = \lim_{i \in I} T_i$ soit une limite filtrante
de schémas affines $T_i$ sur $Y$, et supposons que $F_n(T)$ soit non vide.
Nous devons montrer que $F_n(T_i)$ est non vide pour un certain $i$.
Choisissons un diagramme comme dans la
partie (1) du Lemme \ref{lemma-pre-flat-dimension-n}.
Fixons $i \in I$ et choisissons un ouvert affine $W_i \subset V \times_Y T_i$
s'envoyant surjectivement sur $T_i$. Pour $i' \geq i$, soit $W_{i'}$
l'image réciproque de $W_i$ dans $V \times_Y T_{i'}$, et
soit $W \subset V \times_Y T$ l'image réciproque de $W_i$.
Alors $W = \lim_{i' \geq i} W_i$ est une limite filtrante de schémas
affines sur $V$. À nouveau, d'après le
Lemme \ref{lemma-pre-flat-dimension-n},
il suffit de montrer que $F'_n(W_{i'})$ est non vide pour un certain
$i' \geq i$. Or nous savons que $F'_n(W)$ est non vide,
car, par hypothèse, $F_n(T) = F'_n(V \times_Y T)$
est non vide. Nous pouvons donc appliquer
Compléments sur la platitude, Lemme \ref{flat-lemma-flat-dimension-n},
et conclure.
\end{proof}

\begin{lemma}
\label{lemma-localize-flat-dimension-n}
Dans la Situation \ref{situation-flat-dimension-n}.
Soit $h : X' \to X$ un morphisme \'etale.
Posons $\mathcal{F}' = h^*\mathcal{F}$ et $f' = f \circ h$.
Soit $F_n'$ le foncteur (\ref{equation-flat-dimension-n})
associé à $(f' : X' \to Y, \mathcal{F}')$.
Alors $F_n$ est un sous-foncteur de $F_n'$ et, si
$h(X') \supset \text{Ass}_{X/Y}(\mathcal{F})$, alors $F_n = F'_n$.
\end{lemma}

\begin{proof}
Choisissons $U \to X$, $V \to Y$, $U \to V$ comme dans la partie (1) du
Lemme \ref{lemma-pre-flat-dimension-n}. Choisissons un morphisme \'etale
surjectif $U' \to U \times_X X'$, où $U'$ est un schéma.
Nous disposons alors du lemme pour les deux foncteurs
$F_{U, n}$ et $F_{U', n}$ déterminés par $U' \to U$ et $\mathcal{F}|_U$
sur $V$, voir
Compléments sur la platitude, Lemme \ref{flat-lemma-localize-flat-dimension-n}.
D'autre part, le Lemme \ref{lemma-pre-flat-dimension-n}
nous dit que, pour tout $T \to Y$, on a
$F_n(T) = F_{U, n}(V \times_Y T)$
et
$F'_n(T) = F_{U', n}(V \times_Y T)$.
Cela démontre le lemme.
\end{proof}

\begin{theorem}
\label{theorem-flat-dimension-n-representable}
Dans la Situation \ref{situation-flat-dimension-n}.
Supposons en outre que $f$ soit de présentation finie, que
$\mathcal{F}$ soit un $\mathcal{O}_X$-module de présentation finie
et que $\mathcal{F}$ soit pur relativement à $Y$.
Alors $F_n$ est un espace algébrique et
$F_n \to Y$ est un monomorphisme de présentation finie.
\end{theorem}

\begin{proof}
Le foncteur $F_n$ est un faisceau pour la topologie fppf d'après le
Lemme \ref{lemma-flat-dimension-n}.
Puisque $F_n \to Y$ est un monomorphisme de faisceaux sur
$(\Sch/S)_{fppf}$, nous voyons que $\Delta : F_n \to F_n \times F_n$
est l'image réciproque de la diagonale $\Delta_Y : Y \to Y \times_S Y$.
La représentabilité de $\Delta_Y$ implique donc celle
de la diagonale de $F_n$. Il suffit par conséquent de démontrer
qu'il existe un schéma $W$ sur $S$ et un morphisme \'etale surjectif
$W \to F_n$.

\medskip\noindent
Pour construire $W \to F_n$, choisissons un recouvrement \'etale $\{Y_i \to Y\}$
où $Y_i$ est un schéma. Soit $X_i = X \times_Y Y_i$ et soit
$\mathcal{F}_i$ l'image réciproque de $\mathcal{F}$ sur $X_i$.
Alors $\mathcal{F}_i$ est pur relativement à $Y_i$, soit par définition,
soit d'après le Lemme \ref{lemma-quasi-finite-base-change}.
Les autres hypothèses du théorème sont également préservées.
Enfin, la restriction de $F_n$ à $Y_i$ est
le foncteur $F_n$ correspondant à $X_i \to Y_i$ et $\mathcal{F}_i$.
Il suffit donc de montrer ceci : étant donnés $\mathcal{F}$ et $f : X \to Y$
comme dans l'énoncé du théorème, avec $Y$ schéma, le
foncteur $F_n$ est représentable par un schéma $Z_n$ et
$Z_n \to Y$ est un monomorphisme de présentation finie.

\medskip\noindent
Observons qu'un monomorphisme de présentation finie est
séparé et quasi-fini (Morphismes, Lemme
\ref{morphisms-lemma-monomorphism-loc-finite-type-loc-quasi-finite}).
Ainsi, en combinant
Descente, Lemme \ref{descent-lemma-descent-data-sheaves},
Compléments sur les morphismes, Lemme
\ref{more-morphisms-lemma-separated-locally-quasi-finite-morphisms-fppf-descend}
et
Descente, Lemmes \ref{descent-lemma-descending-property-monomorphism} et
\ref{descent-lemma-descending-property-finite-presentation}
nous voyons que la question est locale pour la topologie \'etale sur $Y$.

\medskip\noindent
En particulier, la situation est locale pour la topologie de Zariski sur $Y$,
et nous pouvons supposer $Y$ affine. Dans ce cas, la dimension des
fibres de $f$ est majorée ; $F_n$ est donc représentable
pour $n$ assez grand. Nous pouvons ainsi raisonner par récurrence descendante sur $n$.
Supposons que nous sachions $F_{n + 1}$ représentable par un monomorphisme
$Z_{n + 1} \to Y$ de présentation finie. Considérons le changement de base
$X_{n + 1} = Z_{n + 1} \times_Y X$ et l'image réciproque $\mathcal{F}_{n + 1}$
de $\mathcal{F}$ sur $X_{n + 1}$. Le morphisme $Z_{n + 1} \to Y$ est
quasi-fini, car c'est un monomorphisme de présentation finie ; le
Lemme \ref{lemma-quasi-finite-base-change}
implique donc que $\mathcal{F}_{n + 1}$ est pur relativement à $Z_{n + 1}$.
Puisque $F_n$ est un sous-foncteur de $F_{n + 1}$, nous concluons que, pour
démontrer le résultat pour $F_n$, il suffit de le démontrer pour le
foncteur correspondant à la situation
$\mathcal{F}_{n + 1}/X_{n + 1}/Z_{n + 1}$.
Nous sommes ainsi ramenés à démontrer le résultat pour $F_n$ dans le cas
$Y_{n + 1} = Y$, c'est-à-dire que nous pouvons supposer $\mathcal{F}$ plat
en dimensions $\geq n + 1$ sur $Y$.

\medskip\noindent
Fixons $n$ et supposons que $\mathcal{F}$ soit plat en dimensions $\geq n + 1$
sur le schéma affine $Y$.
Pour achever la démonstration, nous devons montrer que $F_n$ est représentable
par un monomorphisme $Z_n \to S$ de présentation finie.
Puisque la question est locale pour la topologie \'etale sur $Y$, il suffit de
montrer que, pour tout $y \in Y$, il existe un voisinage \'etale
$(Y', y') \to (Y, y)$ tel que le résultat soit vrai après changement de base à $Y'$.
Ainsi, d'après le
Lemme \ref{lemma-existence-complete},
nous pouvons supposer qu'il existe des morphismes \'etales $h_j : W_j \to X$,
$j = 1, \ldots, m$, tels que, pour chaque $j$, il existe un dévissage
complet de $\mathcal{F}_j/W_j/Y$ au-dessus de $y$,
où $\mathcal{F}_j$ est l'image réciproque de $\mathcal{F}$ sur $W_j$,
et tels que $|X_y| \subset \bigcup h_j(W_j)$. Puisque
$h_j$ est \'etale, le
Lemme \ref{lemma-pre-flat-dimension-n}
montre que les faisceaux $\mathcal{F}_j$ sont encore plats
en dimensions $\geq n + 1$ sur $Y$.
Posons $W = \bigcup h_j(W_j)$, qui est un ouvert quasi-compact de $X$.
Comme $\mathcal{F}$ est pur le long de $X_y$, nous voyons que
$$
E = \{t \in |Y| : \text{Ass}_{X_t}(\mathcal{F}_t) \subset W \}.
$$
contient toutes les générisations de $y$. D'après
Diviseurs sur les espaces,
Lemme \ref{spaces-divisors-lemma-relative-assassin-constructible},
$E$ est une partie constructible de $Y$. Nous avons vu que
$\Spec(\mathcal{O}_{Y, y}) \subset E$. D'après
Morphismes, Lemme \ref{morphisms-lemma-constructible-containing-open},
nous voyons que $E$ contient un voisinage ouvert de $y$.
Ainsi, après avoir restreint $Y$, nous pouvons supposer que $E = Y$.
Il résulte du
Lemme \ref{lemma-localize-flat-dimension-n}
qu'il suffit de démontrer le théorème pour le foncteur $F_n$ associé à
$X = \coprod W_j$ et $\mathcal{F} = \coprod \mathcal{F}_j$.
Si $F_{j, n}$ désigne le foncteur associé à $W_j \to Y$ et au faisceau
$\mathcal{F}_j$, alors $F_n = \prod F_{j, n}$. Il suffit donc
de démontrer que chaque $F_{j, n}$ est représentable par un monomorphisme
$Z_{j, n} \to Y$ de présentation finie, puisque alors
$$
Z_n = Z_{1, n} \times_Y \ldots \times_Y Z_{m, n}
$$
Nous avons ainsi ramené le théorème au cas particulier traité dans
Compléments sur la platitude, Lemme \ref{flat-lemma-flat-dimension-n-representable}.
\end{proof}

\noindent
Nous obtenons enfin le résultat souhaité.

\begin{lemma}
\label{lemma-when-universal-flattening}
Soit $S$ un schéma.
Soit $f : X \to Y$ un morphisme d'espaces algébriques sur $S$.
Soit $\mathcal{F}$ un $\mathcal{O}_X$-module quasi-cohérent.
\begin{enumerate}
\item Si $f$ est de présentation finie, si $\mathcal{F}$ est un
$\mathcal{O}_X$-module de présentation finie et si $\mathcal{F}$ est
pur relativement à $Y$, alors il existe un aplatissement universel
$Y' \to Y$ de $\mathcal{F}$. De plus, $Y' \to Y$ est un monomorphisme
de présentation finie.
\item Si $f$ est de présentation finie et si $X$ est pur relativement à $Y$,
alors il existe un aplatissement universel $Y' \to Y$ de $X$.
De plus, $Y' \to Y$ est un monomorphisme de présentation finie.
\item Si $f$ est propre et de présentation finie et si $\mathcal{F}$ est un
$\mathcal{O}_X$-module de présentation finie, alors il existe un
aplatissement universel $Y' \to Y$ de $\mathcal{F}$. De plus, $Y' \to Y$ est
un monomorphisme de présentation finie.
\item Si $f$ est propre et de présentation finie,
alors il existe un aplatissement universel $Y' \to Y$ de $X$.
\end{enumerate}
\end{lemma}

\begin{proof}
Ces assertions résultent immédiatement du
Théorème \ref{theorem-flat-dimension-n-representable}
appliqué à $F_0 = F_{flat}$,
et du fait que, si $f$ est propre, alors $\mathcal{F}$ est automatiquement
pur sur la base, voir
Lemme \ref{lemma-proper-pure}.
\end{proof}






\section{Théorème d'existence de Grothendieck}
\label{section-existence}

\noindent
Cette section est l'analogue de
Compléments sur la platitude, Section \ref{flat-section-existence},
et poursuit l'étude menée dans
Compléments sur les morphismes d'espaces, Section
\ref{spaces-more-morphisms-section-existence-proper}.
Nous travaillerons dans la situation suivante.

\begin{situation}
\label{situation-existence}
Nous avons ici un système projectif d'anneaux $(A_n)$ dont les applications
de transition sont surjectives et ont des noyaux localement nilpotents. Posons $A = \lim A_n$. Soit
$X$ un espace algébrique séparé et de présentation finie sur $A$.
Posons $X_n = X \times_{\Spec(A)} \Spec(A_n)$ et considérons-le
comme un sous-espace fermé de $X$. Supposons en outre donné un système
$(\mathcal{F}_n, \varphi_n)$, où $\mathcal{F}_n$ est un
$\mathcal{O}_{X_n}$-module de présentation finie, plat sur $A_n$, de support propre sur $A_n$,
et où
$$
\varphi_n :
\mathcal{F}_n \otimes_{\mathcal{O}_{X_n}} \mathcal{O}_{X_{n - 1}}
\longrightarrow
\mathcal{F}_{n - 1}
$$
est un isomorphisme (notation utilisant l'équivalence de
Morphismes d'espaces, Lemme \ref{spaces-morphisms-lemma-i-star-equivalence}).
\end{situation}

\noindent
Notre but est de déterminer si l'on peut trouver un faisceau quasi-cohérent $\mathcal{F}$
sur $X$ tel que
$\mathcal{F}_n = \mathcal{F} \otimes_{\mathcal{O}_X} \mathcal{O}_{X_n}$
pour tout $n$.

\begin{lemma}
\label{lemma-compute-what-it-should-be}
Dans la Situation \ref{situation-existence}, considérons
$$
K = R\lim_{D_\QCoh(\mathcal{O}_X)}(\mathcal{F}_n) =
DQ_X(R\lim_{D(\mathcal{O}_X)}\mathcal{F}_n)
$$
Alors $K$ appartient à $D^b_{\QCoh}(\mathcal{O}_X)$ et, en fait,
$K$ n'a de faisceaux de cohomologie non nuls qu'en degrés $\geq 0$.
\end{lemma}

\begin{proof}
C'est un cas particulier de
Catégories dérivées des espaces, Exemple
\ref{spaces-perfect-example-inverse-limit-quasi-coherent}.
\end{proof}

\begin{lemma}
\label{lemma-compute-against-perfect}
Dans la Situation \ref{situation-existence}, soit $K$ comme dans le
Lemme \ref{lemma-compute-what-it-should-be}. Pour tout objet parfait
$E$ de $D(\mathcal{O}_X)$, on a :
\begin{enumerate}
\item $M = R\Gamma(X, K \otimes^\mathbf{L} E)$ est un objet parfait de $D(A)$,
et il existe un isomorphisme canonique
$R\Gamma(X_n, \mathcal{F}_n \otimes^\mathbf{L} E|_{X_n}) =
M \otimes_A^\mathbf{L} A_n$
dans $D(A_n)$,
\item $N = R\Hom_X(E, K)$ est un objet parfait de $D(A)$,
et il existe un isomorphisme canonique
$R\Hom_{X_n}(E|_{X_n}, \mathcal{F}_n) = N \otimes_A^\mathbf{L} A_n$
dans $D(A_n)$.
\end{enumerate}
Dans les deux assertions, $E|_{X_n}$ désigne l'image réciproque dérivée
de $E$ sur $X_n$.
\end{lemma}

\begin{proof}
Preuve de (2). Écrivons $E_n = E|_{X_n}$ et
$N_n = R\Hom_{X_n}(E_n, \mathcal{F}_n)$.
Rappelons que $R\Hom_{X_n}(-, -)$ est égal à
$R\Gamma(X_n, R\SheafHom(-, -))$, voir
Cohomologie sur les sites, Section \ref{sites-cohomology-section-global-RHom}.
Ainsi, d'après Catégories dérivées des espaces, Lemme
\ref{spaces-perfect-lemma-base-change-RHom-perfect}
le complexe $N_n$ est un objet parfait de $D(A_n)$
dont la formation commute au changement de base. Ainsi, les applications
$N_n \otimes_{A_n}^\mathbf{L} A_{n - 1} \to N_{n - 1}$
provenant de $\varphi_n$ sont des isomorphismes.
D'après Compléments d'algèbre, Lemme \ref{more-algebra-lemma-Rlim-perfect-gives-perfect},
$R\lim N_n$ est parfait et
son changement de base à $A_n$ redonne $N_n$.
D'autre part, le foncteur exact
$R\Hom_X(E, -) : D_\QCoh(\mathcal{O}_X) \to D(A)$
de catégories triangulées commute aux produits,
donc aux limites dérivées ; par conséquent,
$$
R\Hom_X(E, K) =
R\lim R\Hom_X(E, \mathcal{F}_n) =
R\lim R\Hom_X(E_n, \mathcal{F}_n) =
R\lim N_n
$$
Cela démontre (2). Pour obtenir (1), on le reformule sous la forme (2)
à l'aide de Cohomologie sur les sites, Lemme
\ref{sites-cohomology-lemma-dual-perfect-complex}.
\end{proof}

\begin{lemma}
\label{lemma-relative-pseudo-coherence}
Dans la Situation \ref{situation-existence}, soit $K$ comme dans le
Lemme \ref{lemma-compute-what-it-should-be}. Alors $K$
est pseudo-cohérent relativement à $A$.
\end{lemma}

\begin{proof}
En combinant le Lemme \ref{lemma-compute-against-perfect} et
Catégories dérivées des espaces, Lemme \ref{spaces-perfect-lemma-perfect-enough},
nous voyons que $R\Gamma(X, K \otimes^\mathbf{L} E)$
est pseudo-cohérent dans $D(A)$ pour tout objet pseudo-cohérent
$E$ de $D(\mathcal{O}_X)$. Le lemme résulte donc de
Compléments sur les morphismes d'espaces, Lemme
\ref{spaces-more-morphisms-lemma-characterize-pseudo-coherent}.
\end{proof}

\begin{lemma}
\label{lemma-compute-over-affine}
Dans la Situation \ref{situation-existence}, soit $K$ comme dans le
Lemme \ref{lemma-compute-what-it-should-be}. Pour tout
morphisme \'etale $U \to X$, avec $U$ quasi-compact et quasi-séparé, on a
$$
R\Gamma(U, K) \otimes_A^\mathbf{L} A_n =
R\Gamma(U_n, \mathcal{F}_n)
$$
dans $D(A_n)$, où $U_n = U \times_X X_n$.
\end{lemma}

\begin{proof}
Fixons $n$. D'après Catégories dérivées des espaces, Lemme
\ref{spaces-perfect-lemma-computing-sections-as-colim}
il existe un système de complexes parfaits $E_m$
sur $X$ tel que
$R\Gamma(U, K) = \text{hocolim} R\Gamma(X, K \otimes^\mathbf{L} E_m)$.
En fait, cette formule vaut non seulement pour $K$, mais pour tout objet de
$D_\QCoh(\mathcal{O}_X)$.
En l'appliquant à $\mathcal{F}_n$,
nous obtenons
\begin{align*}
R\Gamma(U_n, \mathcal{F}_n)
& =
R\Gamma(U, \mathcal{F}_n) \\
& =
\text{hocolim}_m R\Gamma(X, \mathcal{F}_n \otimes^\mathbf{L} E_m) \\
& =
\text{hocolim}_m R\Gamma(X_n, \mathcal{F}_n \otimes^\mathbf{L} E_m|_{X_n})
\end{align*}
En utilisant le Lemme \ref{lemma-compute-against-perfect}
et le fait que $- \otimes_A^\mathbf{L} A_n$
commute aux colimites homotopiques, nous obtenons le résultat.
\end{proof}

\begin{lemma}
\label{lemma-finitely-presented}
Dans la Situation \ref{situation-existence}, soit $K$ comme dans le
Lemme \ref{lemma-compute-what-it-should-be}.
Notons $X_0 \subset |X|$ la partie fermée
formée des points situés au-dessus de la partie fermée
$\Spec(A_1) = \Spec(A_2) = \ldots$ de $\Spec(A)$.
Il existe un sous-espace ouvert $W \subset X$ contenant $X_0$
tel que
\begin{enumerate}
\item $H^i(K)|_W$ soit nul sauf si $i = 0$,
\item $\mathcal{F} = H^0(K)|_W$ soit de présentation finie, et
\item $\mathcal{F}_n = \mathcal{F} \otimes_{\mathcal{O}_X} \mathcal{O}_{X_n}$.
\end{enumerate}
\end{lemma}

\begin{proof}
Fixons $n \geq 1$. Par construction, il existe une application canonique
$K \to \mathcal{F}_n$ dans $D_\QCoh(\mathcal{O}_X)$,
et donc une application canonique $H^0(K) \to \mathcal{F}_n$
de faisceaux quasi-cohérents. Cela précise le sens de la partie (3).

\medskip\noindent
Soit $x \in X_0$ un point. Nous allons trouver un voisinage ouvert $W$
de $x$ où (1), (2) et (3) sont vraies. Puisque $X_0$ est quasi-compact,
cela démontrera le lemme. Soit $U \to X$ un morphisme \'etale,
avec $U$ affine, et soit $u \in U$ un point s'envoyant sur $x$. Puisque $|U| \to |X|$
est ouvert, il suffit de trouver un voisinage ouvert de $u$ dans $U$
où (1), (2) et (3) sont vraies. Écrivons $U = \Spec(B)$.
Choisissons une surjection $P \to B$, avec $P$ lisse sur $A$.
D'après le Lemme \ref{lemma-relative-pseudo-coherence}
et la définition de la pseudo-cohérence relative,
il existe un complexe borné supérieurement $F^\bullet$
de $P$-modules libres de rang fini représentant
$Ri_*K$, où $i : U \to \Spec(P)$ est l'immersion
fermée induite par la présentation.
Soit $M_n$ le $B$-module correspondant à $\mathcal{F}_n|_U$.
D'après le Lemme \ref{lemma-compute-over-affine},
$$
H^i(F^\bullet \otimes_A A_n) =
\left\{
\begin{matrix}
0 & \text{si} & i \not = 0 \\
M_n & \text{si} & i = 0
\end{matrix}
\right.
$$
Soit $i$ l'indice maximal tel que $F^i$ soit non nul.
Si $i \leq 0$, alors (1), (2) et (3) sont vraies.
Sinon, $i > 0$, et nous voyons que le rang de l'application
$$
F^{i - 1} \to F^i
$$
au point $u$ est maximal. Il est donc maximal dans un voisinage ouvert
de $u$ dans $\Spec(P)$. Ainsi, après avoir remplacé
$P$ par une localisation principale, nous pouvons supposer que l'application affichée
est surjective. Puisque $F^i$ est libre de rang fini, nous pouvons choisir
une décomposition $F^{i - 1} = F' \oplus F^i$. Nous pouvons alors
remplacer $F^\bullet$ par le complexe
$$
\ldots \to F^{i - 2} \to F' \to 0 \to \ldots
$$
et conclure par récurrence sur $i$.
\end{proof}

\begin{lemma}
\label{lemma-proper-support}
Dans la Situation \ref{situation-existence}, soit $K$ comme dans le
Lemme \ref{lemma-compute-what-it-should-be}. Soit $W \subset X$
comme dans le Lemme \ref{lemma-finitely-presented}.
Posons $\mathcal{F} = H^0(K)|_W$. Alors, quitte à restreindre l'ouvert $W$,
le support de $\mathcal{F}$ est propre sur $A$.
\end{lemma}

\begin{proof}
Fixons $n \geq 1$. Soit $I_n = \Ker(A \to A_n)$.
D'après Compléments d'algèbre, Lemme \ref{more-algebra-lemma-limit-henselian},
la paire $(A, I_n)$ est hensélienne.
Soit $Z \subset W$ le support schématique de $\mathcal{F}$.
C'est un sous-espace fermé, puisque $\mathcal{F}$ est de présentation finie.
D'après la partie (3) du Lemme \ref{lemma-finitely-presented},
$Z \times_{\Spec(A)} \Spec(A_n)$
est égal au support de $\mathcal{F}_n$, et est donc
propre sur $\Spec(A/I)$.
D'après Compléments sur les morphismes d'espaces, Lemme
\ref{spaces-more-morphisms-lemma-split-off-proper-part-henselian}
nous pouvons écrire $Z = Z_1 \amalg Z_2$, avec $Z_1, Z_2$ ouverts et
fermés dans $Z$, avec $Z_1$ propre
sur $A$, et avec $Z_1 \times_{\Spec(A)} \Spec(A/I_n)$
égal au support de $\mathcal{F}_n$.
Autrement dit, $|Z_2|$ ne rencontre pas $X_0$.
Ainsi, en remplaçant $W$ par $W \setminus Z_2$, nous obtenons le lemme.
\end{proof}

\begin{theorem}[Théorème d'existence de Grothendieck]
\label{theorem-existence}
Dans la Situation \ref{situation-existence},
il existe un $\mathcal{O}_X$-module de présentation finie
$\mathcal{F}$, plat sur $A$, de support propre sur $A$,
tel que
$\mathcal{F}_n = \mathcal{F} \otimes_{\mathcal{O}_X} \mathcal{O}_{X_n}$
pour tout $n$, de manière compatible avec les applications $\varphi_n$.
\end{theorem}

\begin{proof}
Appliquons les Lemmes \ref{lemma-compute-what-it-should-be},
\ref{lemma-compute-against-perfect},
\ref{lemma-relative-pseudo-coherence},
\ref{lemma-compute-over-affine},
\ref{lemma-finitely-presented} et
\ref{lemma-proper-support}
pour obtenir un sous-espace ouvert $W \subset X$ contenant tous les points
au-dessus de $\Spec(A_n)$
et un $\mathcal{O}_W$-module $\mathcal{F}$ de présentation finie,
dont le support est propre sur $A$ et tel que
$\mathcal{F}_n = \mathcal{F} \otimes_{\mathcal{O}_W} \mathcal{O}_{X_n}$
pour tout $n \geq 1$. (Cela a un sens puisque $X_n \subset W$.)
D'après le Lemme \ref{lemma-proper-pure}, $\mathcal{F}$
est universellement pur relativement à $\Spec(A)$.
D'après le Théorème \ref{theorem-flat-dimension-n-representable}
(pour des explications, voir le Lemme \ref{lemma-when-universal-flattening}),
il existe un aplatissement universel $S' \to \Spec(A)$
de $\mathcal{F}$ et, de plus, le morphisme $S' \to \Spec(A)$
est un monomorphisme de présentation finie.
En particulier, $S'$ est un schéma (cela résulte de la démonstration
du théorème, mais aussi a posteriori de
Morphismes d'espaces, Proposition
\ref{spaces-morphisms-proposition-locally-quasi-finite-separated-over-scheme}).
Puisque le changement de base de $\mathcal{F}$ à $\Spec(A_n)$
est $\mathcal{F}_n$, nous voyons que $\Spec(A_n) \to \Spec(A)$
se factorise (uniquement) par $S'$ pour tout $n$.
D'après Compléments sur la platitude, Lemme \ref{flat-lemma-monomorphism-isomorphism},
nous avons $S' = \Spec(A)$.
Cela signifie que $\mathcal{F}$ est plat sur $A$.
Enfin, puisque le support schématique $Z$ de $\mathcal{F}$
est propre sur $\Spec(A)$, le morphisme $Z \to X$ est fermé.
Par conséquent, l'image directe $(W \to X)_*\mathcal{F}$ est supportée
sur $W$ et possède toutes les propriétés voulues.
\end{proof}







\section{Théorème d'existence de Grothendieck, bis}
\label{section-existence-derived}

\noindent
Dans cette section, nous démontrons un analogue du théorème d'existence de Grothendieck
dans la catégorie dérivée, en suivant la méthode utilisée à la
Section \ref{section-existence} pour les modules quasi-cohérents.
Cette section est l'analogue, pour les espaces algébriques, de
Compléments sur la platitude, Section \ref{flat-section-existence-derived}.
Le cas classique (pour les espaces algébriques)
est étudié dans Compléments sur les morphismes d'espaces, Section
\ref{spaces-more-morphisms-section-existence-proper}.
Nous travaillerons dans la situation suivante.

\begin{situation}
\label{situation-existence-derived}
Nous avons ici un système projectif d'anneaux $(A_n)$ dont les applications
de transition sont surjectives et ont des noyaux localement nilpotents. Posons $A = \lim A_n$. Soit
$X$ un espace algébrique propre, plat et de présentation finie sur $A$.
Posons $X_n = X \times_{\Spec(A)} \Spec(A_n)$ et considérons-le
comme un sous-espace fermé de $X$. Supposons en outre donné un système
$(K_n, \varphi_n)$, où $K_n$ est un objet pseudo-cohérent de
$D(\mathcal{O}_{X_n})$ et
$$
\varphi_n : K_n \longrightarrow K_{n - 1}
$$
est une application dans $D(\mathcal{O}_{X_n})$ qui induit un isomorphisme
$K_n \otimes_{\mathcal{O}_{X_n}}^\mathbf{L} \mathcal{O}_{X_{n - 1}}
\to K_{n - 1}$ dans $D(\mathcal{O}_{X_{n - 1}})$.
\end{situation}

\noindent
Plus précisément, il faudrait écrire
$\varphi_n : K_n \to Ri_{n - 1, *}K_{n - 1}$
où $i_{n - 1} : X_{n - 1} \to X_n$ est le morphisme d'inclusion ;
avec cette notation, la condition est que l'application adjointe
$Li_{n - 1}^*K_n \to K_{n - 1}$ soit un isomorphisme.
Notre but est de trouver un $K \in D(\mathcal{O}_X)$ pseudo-cohérent
tel que $K_n = K \otimes_{\mathcal{O}_X}^\mathbf{L} \mathcal{O}_{X_n}$
pour tout $n$ (avec le même abus de notation).

\begin{lemma}
\label{lemma-compute-what-it-should-be-derived}
Dans la Situation \ref{situation-existence-derived}, considérons
$$
K = R\lim_{D_\QCoh(\mathcal{O}_X)}(K_n) =
DQ_X(R\lim_{D(\mathcal{O}_X)} K_n)
$$
Alors $K$ appartient à $D^-_{\QCoh}(\mathcal{O}_X)$.
\end{lemma}

\begin{proof}
Le foncteur $DQ_X$ existe car $X$ est quasi-compact et
quasi-séparé, voir Catégories dérivées des espaces, Lemme
\ref{spaces-perfect-lemma-better-coherator}.
Puisque $DQ_X$ est un adjoint à droite, il commute aux produits,
et donc aux limites dérivées. On obtient ainsi l'égalité
de l'énoncé du lemme.

\medskip\noindent
D'après Catégories dérivées des espaces,
Lemme \ref{spaces-perfect-lemma-boundedness-better-coherator},
le foncteur $DQ_X$ est de dimension cohomologique bornée.
Il suffit donc de montrer que $R\lim K_n \in D^-(\mathcal{O}_X)$.
Pour cela, soit $U \to X$ un morphisme \'etale, avec $U$ affine.
Il existe alors une suite exacte canonique
$$
0 \to
R^1\lim H^{m - 1}(U, K_n) \to H^m(U, R\lim K_n) \to
\lim H^m(U, K_n) \to 0
$$
d'après Cohomologie sur les sites, Lemme
\ref{sites-cohomology-lemma-RGamma-commutes-with-Rlim}.
Puisque $U$ est affine et $K_n$ pseudo-cohérent (et possède donc des
faisceaux de cohomologie quasi-cohérents d'après
Catégories dérivées des espaces, Lemme \ref{spaces-perfect-lemma-pseudo-coherent}),
nous avons $H^m(U, K_n) = H^m(K_n)(U)$ d'après
Catégories dérivées des schémas, Lemme \ref{perfect-lemma-affine-compare-bounded}.
Il suffit donc de montrer que les $K_n$ sont
bornés supérieurement indépendamment de $n$.

\medskip\noindent
Puisque $K_n$ est pseudo-cohérent, on a $K_n \in D^-(\mathcal{O}_{X_n})$.
Supposons que $a_n$ soit maximal tel que $H^{a_n}(K_n)$ soit non nul.
Bien entendu, $a_1 \leq a_2 \leq a_3 \leq \ldots$.
Remarquons que $H^{a_n}(K_n)$ est un
$\mathcal{O}_{X_n}$-module de présentation finie
(Cohomologie sur les sites, Lemme
\ref{sites-cohomology-lemma-finite-cohomology}).
Nous avons $H^{a_n}(K_{n - 1}) =
H^{a_n}(K_n) \otimes_{\mathcal{O}_{X_n}} \mathcal{O}_{X_{n - 1}}$.
Puisque $X_{n - 1} \to X_n$ est un épaississement, le
lemme de Nakayama (Algèbre, Lemme \ref{algebra-lemma-NAK}) montre que, si
$H^{a_n}(K_n) \otimes_{\mathcal{O}_{X_n}} \mathcal{O}_{X_{n - 1}}$
est nul, alors $H^{a_n}(K_n)$ l'est aussi (on peut par exemple vérifier
sur les germes ; ce petit détail est omis).
Ainsi $a_{n - 1} = a_n$ pour tout $n$, ce qui permet de conclure.
\end{proof}

\begin{lemma}
\label{lemma-compute-against-perfect-derived}
Dans la Situation \ref{situation-existence-derived}, soit $K$ comme dans le
Lemme \ref{lemma-compute-what-it-should-be-derived}. Pour tout objet parfait
$E$ de $D(\mathcal{O}_X)$, la cohomologie
$$
M = R\Gamma(X, K \otimes^\mathbf{L} E)
$$
est un objet pseudo-cohérent de $D(A)$ et il existe un isomorphisme canonique
$$
R\Gamma(X_n, K_n \otimes^\mathbf{L} E|_{X_n}) = M \otimes_A^\mathbf{L} A_n
$$
dans $D(A_n)$. Ici, $E|_{X_n}$ désigne l'image réciproque dérivée de $E$ sur $X_n$.
\end{lemma}

\begin{proof}
Écrivons $E_n = E|_{X_n}$ et
$M_n = R\Gamma(X_n, K_n \otimes^\mathbf{L} E|_{X_n})$.
D'après Catégories dérivées des espaces, Lemme
\ref{spaces-perfect-lemma-flat-proper-pseudo-coherent-direct-image-general}
nous voyons que $M_n$ est un objet pseudo-cohérent de $D(A_n)$
dont la formation commute au changement de base. Ainsi, les applications
$M_n \otimes_{A_n}^\mathbf{L} A_{n - 1} \to M_{n - 1}$
provenant de $\varphi_n$ sont des isomorphismes. D'après
Compléments d'algèbre, Lemme
\ref{more-algebra-lemma-Rlim-pseudo-coherent-gives-pseudo-coherent}
$R\lim M_n$ est pseudo-cohérent et
son changement de base à $A_n$ redonne $M_n$.
D'autre part, le foncteur exact
$R\Gamma(X, -) : D_\QCoh(\mathcal{O}_X) \to D(A)$
de catégories triangulées commute aux produits,
donc aux limites dérivées ; par conséquent,
$$
R\Gamma(X, E \otimes^\mathbf{L} K) =
R\lim R\Gamma(X,  E \otimes^\mathbf{L} K_n) =
R\lim R\Gamma(X_n, E_n \otimes^\mathbf{L} K_n) =
R\lim M_n
$$
comme souhaité.
\end{proof}

\begin{lemma}
\label{lemma-relative-pseudo-coherence-derived}
Dans la Situation \ref{situation-existence-derived}, soit $K$ comme dans le
Lemme \ref{lemma-compute-what-it-should-be-derived}. Alors $K$
est pseudo-cohérent sur $X$.
\end{lemma}

\begin{proof}
En combinant le Lemme \ref{lemma-compute-against-perfect-derived} et
Catégories dérivées des espaces, Lemme
\ref{spaces-perfect-lemma-perfect-enough}
nous voyons que $R\Gamma(X, K \otimes^\mathbf{L} E)$
est pseudo-cohérent dans $D(A)$ pour tout objet pseudo-cohérent
$E$ de $D(\mathcal{O}_X)$. Il résulte donc de
Compléments sur les morphismes d'espaces, Lemme
\ref{spaces-more-morphisms-lemma-characterize-pseudo-coherent}
que $K$ est pseudo-cohérent relativement à $A$.
Puisque $X$ est plat et de présentation finie
sur $A$, cela équivaut à être pseudo-cohérent sur $X$, voir
Compléments sur les morphismes d'espaces, Lemme
\ref{spaces-more-morphisms-lemma-relative-pseudo-coherent-is-moot}.
\end{proof}

\begin{lemma}
\label{lemma-compute-over-affine-derived}
Dans la Situation \ref{situation-existence-derived}, soit $K$ comme dans le
Lemme \ref{lemma-compute-what-it-should-be-derived}. Pour tout
morphisme \'etale $U \to X$, avec $U$ quasi-compact et quasi-séparé, on a
$$
R\Gamma(U, K) \otimes_A^\mathbf{L} A_n =
R\Gamma(U_n, K_n)
$$
dans $D(A_n)$, où $U_n = U \times_X X_n$.
\end{lemma}

\begin{proof}
Fixons $n$. D'après Catégories dérivées des espaces, Lemme
\ref{spaces-perfect-lemma-computing-sections-as-colim}
il existe un système de complexes parfaits $E_m$
sur $X$ tel que
$R\Gamma(U, K) = \text{hocolim} R\Gamma(X, K \otimes^\mathbf{L} E_m)$.
En fait, cette formule vaut non seulement pour $K$, mais pour tout objet de
$D_\QCoh(\mathcal{O}_X)$.
En l'appliquant à $K_n$,
nous obtenons
\begin{align*}
R\Gamma(U_n, K_n)
& =
R\Gamma(U, K_n) \\
& =
\text{hocolim}_m R\Gamma(X, K_n \otimes^\mathbf{L} E_m) \\
& =
\text{hocolim}_m R\Gamma(X_n, K_n \otimes^\mathbf{L} E_m|_{X_n})
\end{align*}
En utilisant le Lemme \ref{lemma-compute-against-perfect-derived}
et le fait que $- \otimes_A^\mathbf{L} A_n$
commute aux colimites homotopiques, nous obtenons le résultat.
\end{proof}

\begin{theorem}[Théorème d'existence de Grothendieck dérivé]
\label{theorem-existence-derived}
Dans la Situation \ref{situation-existence-derived},
il existe un $K$ pseudo-cohérent dans $D(\mathcal{O}_X)$
tel que $K_n = K \otimes_{\mathcal{O}_X}^\mathbf{L} \mathcal{O}_{X_n}$
pour tout $n$, de manière compatible avec les applications $\varphi_n$.
\end{theorem}

\begin{proof}
Appliquons les Lemmes \ref{lemma-compute-what-it-should-be-derived},
\ref{lemma-compute-against-perfect-derived},
\ref{lemma-relative-pseudo-coherence-derived}
pour obtenir un objet pseudo-cohérent $K$ de $D(\mathcal{O}_X)$.
En choisissant $U$ affine dans le Lemme
\ref{lemma-compute-over-affine-derived}
on voit immédiatement que $K$ se restreint en $K_n$ sur $X_n$.
\end{proof}

\begin{remark}
\label{remark-correct-generality}
Le résultat de cette section peut être généralisé. Il est probablement vrai
si l'on suppose seulement que $X \to \Spec(A)$ est séparé et de présentation finie,
et que $K_n$ est pseudo-cohérent relativement à $A_n$ et supporté par une partie
fermée de $X_n$ propre sur $A_n$. On obtiendra alors un $K$
pseudo-cohérent relativement à $A$ et supporté par une partie fermée
propre sur $A$. Si cela nous est un jour nécessaire, nous
formulerons ici un énoncé précis et le démontrerons.
\end{remark}










\begin{multicols}{2}[\section{Autres chapitres}]
\noindent
Préliminaires
\begin{enumerate}
\item \hyperref[introduction-section-phantom]{Introduction}
\item \hyperref[conventions-section-phantom]{Conventions}
\item \hyperref[sets-section-phantom]{Théorie des ensembles}
\item \hyperref[categories-section-phantom]{Catégories}
\item \hyperref[topology-section-phantom]{Topologie}
\item \hyperref[sheaves-section-phantom]{Faisceaux sur les espaces}
\item \hyperref[sites-section-phantom]{Sites et faisceaux}
\item \hyperref[stacks-section-phantom]{Champs}
\item \hyperref[fields-section-phantom]{Corps}
\item \hyperref[algebra-section-phantom]{Algèbre commutative}
\item \hyperref[brauer-section-phantom]{Groupes de Brauer}
\item \hyperref[homology-section-phantom]{Algèbre homologique}
\item \hyperref[derived-section-phantom]{Catégories dérivées}
\item \hyperref[simplicial-section-phantom]{Méthodes simpliciales}
\item \hyperref[more-algebra-section-phantom]{Compléments d'algèbre}
\item \hyperref[smoothing-section-phantom]{Lissage des morphismes d'anneaux}
\item \hyperref[modules-section-phantom]{Faisceaux de modules}
\item \hyperref[sites-modules-section-phantom]{Modules sur les sites}
\item \hyperref[injectives-section-phantom]{Objets injectifs}
\item \hyperref[cohomology-section-phantom]{Cohomologie des faisceaux}
\item \hyperref[sites-cohomology-section-phantom]{Cohomologie sur les sites}
\item \hyperref[dga-section-phantom]{Algèbre différentielle graduée}
\item \hyperref[dpa-section-phantom]{Algèbre à puissances divisées}
\item \hyperref[sdga-section-phantom]{Faisceaux différentiels gradués}
\item \hyperref[hypercovering-section-phantom]{Hyperrecouvrements}
\end{enumerate}
Schémas
\begin{enumerate}
\setcounter{enumi}{25}
\item \hyperref[schemes-section-phantom]{Schémas}
\item \hyperref[constructions-section-phantom]{Constructions de schémas}
\item \hyperref[properties-section-phantom]{Propriétés des schémas}
\item \hyperref[morphisms-section-phantom]{Morphismes de schémas}
\item \hyperref[coherent-section-phantom]{Cohomologie des schémas}
\item \hyperref[divisors-section-phantom]{Diviseurs}
\item \hyperref[limits-section-phantom]{Limites de schémas}
\item \hyperref[varieties-section-phantom]{Variétés}
\item \hyperref[topologies-section-phantom]{Topologies sur les schémas}
\item \hyperref[descent-section-phantom]{Descente}
\item \hyperref[perfect-section-phantom]{Catégories dérivées de schémas}
\item \hyperref[more-morphisms-section-phantom]{Compléments sur les morphismes}
\item \hyperref[flat-section-phantom]{Compléments sur la platitude}
\item \hyperref[groupoids-section-phantom]{Schémas en groupoïdes}
\item \hyperref[more-groupoids-section-phantom]{Compléments sur les schémas en groupoïdes}
\item \hyperref[etale-section-phantom]{Morphismes étales de schémas}
\end{enumerate}
Thèmes de théorie des schémas
\begin{enumerate}
\setcounter{enumi}{41}
\item \hyperref[chow-section-phantom]{Homologie de Chow}
\item \hyperref[intersection-section-phantom]{Théorie de l'intersection}
\item \hyperref[pic-section-phantom]{Schémas de Picard des courbes}
\item \hyperref[weil-section-phantom]{Théories cohomologiques de Weil}
\item \hyperref[adequate-section-phantom]{Modules adéquats}
\item \hyperref[dualizing-section-phantom]{Complexes dualisants}
\item \hyperref[duality-section-phantom]{Dualité pour les schémas}
\item \hyperref[discriminant-section-phantom]{Discriminants et différentes}
\item \hyperref[derham-section-phantom]{Cohomologie de de Rham}
\item \hyperref[local-cohomology-section-phantom]{Cohomologie locale}
\item \hyperref[algebraization-section-phantom]{Géométrie algébrique et formelle}
\item \hyperref[curves-section-phantom]{Courbes algébriques}
\item \hyperref[resolve-section-phantom]{Résolution des surfaces}
\item \hyperref[models-section-phantom]{Réduction semi-stable}
\item \hyperref[functors-section-phantom]{Foncteurs et morphismes}
\item \hyperref[equiv-section-phantom]{Catégories dérivées de variétés}
\item \hyperref[pione-section-phantom]{Groupes fondamentaux de schémas}
\item \hyperref[etale-cohomology-section-phantom]{Cohomologie étale}
\item \hyperref[crystalline-section-phantom]{Cohomologie cristalline}
\item \hyperref[proetale-section-phantom]{Cohomologie pro-étale}
\item \hyperref[relative-cycles-section-phantom]{Cycles relatifs}
\item \hyperref[more-etale-section-phantom]{Compléments de cohomologie étale}
\item \hyperref[trace-section-phantom]{La formule des traces}
\end{enumerate}
Espaces algébriques
\begin{enumerate}
\setcounter{enumi}{64}
\item \hyperref[spaces-section-phantom]{Espaces algébriques}
\item \hyperref[spaces-properties-section-phantom]{Propriétés des espaces algébriques}
\item \hyperref[spaces-morphisms-section-phantom]{Morphismes d'espaces algébriques}
\item \hyperref[decent-spaces-section-phantom]{Espaces algébriques décents}
\item \hyperref[spaces-cohomology-section-phantom]{Cohomologie des espaces algébriques}
\item \hyperref[spaces-limits-section-phantom]{Limites d'espaces algébriques}
\item \hyperref[spaces-divisors-section-phantom]{Diviseurs sur les espaces algébriques}
\item \hyperref[spaces-over-fields-section-phantom]{Espaces algébriques sur des corps}
\item \hyperref[spaces-topologies-section-phantom]{Topologies sur les espaces algébriques}
\item \hyperref[spaces-descent-section-phantom]{Descente et espaces algébriques}
\item \hyperref[spaces-perfect-section-phantom]{Catégories dérivées d'espaces}
\item \hyperref[spaces-more-morphisms-section-phantom]{Compléments sur les morphismes d'espaces}
\item \hyperref[spaces-flat-section-phantom]{Platitude sur les espaces algébriques}
\item \hyperref[spaces-groupoids-section-phantom]{Groupoïdes dans les espaces algébriques}
\item \hyperref[spaces-more-groupoids-section-phantom]{Compléments sur les groupoïdes dans les espaces}
\item \hyperref[bootstrap-section-phantom]{Amorçage}
\item \hyperref[spaces-pushouts-section-phantom]{Sommes amalgamées d'espaces algébriques}
\end{enumerate}
Thèmes de géométrie
\begin{enumerate}
\setcounter{enumi}{81}
\item \hyperref[spaces-chow-section-phantom]{Groupes de Chow des espaces}
\item \hyperref[groupoids-quotients-section-phantom]{Quotients de groupoïdes}
\item \hyperref[spaces-more-cohomology-section-phantom]{Compléments sur la cohomologie des espaces}
\item \hyperref[spaces-simplicial-section-phantom]{Espaces simpliciaux}
\item \hyperref[spaces-duality-section-phantom]{Dualité pour les espaces}
\item \hyperref[formal-spaces-section-phantom]{Espaces algébriques formels}
\item \hyperref[restricted-section-phantom]{Algébrisation des espaces formels}
\item \hyperref[spaces-resolve-section-phantom]{Résolution des surfaces revisitée}
\end{enumerate}
Théorie des déformations
\begin{enumerate}
\setcounter{enumi}{89}
\item \hyperref[formal-defos-section-phantom]{Théorie formelle des déformations}
\item \hyperref[defos-section-phantom]{Théorie des déformations}
\item \hyperref[cotangent-section-phantom]{Le complexe cotangent}
\item \hyperref[examples-defos-section-phantom]{Problèmes de déformation}
\end{enumerate}
Champs algébriques
\begin{enumerate}
\setcounter{enumi}{93}
\item \hyperref[algebraic-section-phantom]{Champs algébriques}
\item \hyperref[examples-stacks-section-phantom]{Exemples de champs}
\item \hyperref[stacks-sheaves-section-phantom]{Faisceaux sur les champs algébriques}
\item \hyperref[criteria-section-phantom]{Critères de représentabilité}
\item \hyperref[artin-section-phantom]{Axiomes d'Artin}
\item \hyperref[quot-section-phantom]{Espaces de Quot et de Hilbert}
\item \hyperref[stacks-properties-section-phantom]{Propriétés des champs algébriques}
\item \hyperref[stacks-morphisms-section-phantom]{Morphismes de champs algébriques}
\item \hyperref[stacks-limits-section-phantom]{Limites de champs algébriques}
\item \hyperref[stacks-cohomology-section-phantom]{Cohomologie des champs algébriques}
\item \hyperref[stacks-perfect-section-phantom]{Catégories dérivées de champs}
\item \hyperref[stacks-introduction-section-phantom]{Introduction aux champs algébriques}
\item \hyperref[stacks-more-morphisms-section-phantom]{Compléments sur les morphismes de champs}
\item \hyperref[stacks-geometry-section-phantom]{La géométrie des champs}
\end{enumerate}
Thèmes de théorie des espaces de modules
\begin{enumerate}
\setcounter{enumi}{107}
\item \hyperref[moduli-section-phantom]{Champs de modules}
\item \hyperref[moduli-curves-section-phantom]{Espaces de modules de courbes}
\end{enumerate}
Divers
\begin{enumerate}
\setcounter{enumi}{109}
\item \hyperref[examples-section-phantom]{Exemples}
\item \hyperref[exercises-section-phantom]{Exercices}
\item \hyperref[guide-section-phantom]{Guide bibliographique}
\item \hyperref[desirables-section-phantom]{Desiderata}
\item \hyperref[coding-section-phantom]{Style de programmation}
\item \hyperref[obsolete-section-phantom]{Obsolète}
\item \hyperref[fdl-section-phantom]{Licence de documentation libre GNU}
\item \hyperref[index-section-phantom]{Index généré automatiquement}
\end{enumerate}
\end{multicols}



\bibliography{my}
\bibliographystyle{amsalpha}

\end{document}
